% Minimal preamble for Open Logic Project sources (pinned commit 1e960beff9ed7835bf3e3f1335e21af3439cd107).
%
% OLP's own master preamble cannot run on the project host: TeX Live 2025 there lacks
% prooftree.sty, ntheorem.sty, extarrows.sty, mathdots.sty, isomath.sty and subfiles.sty.
% Following the project's own rule ("if the author preamble needs packages the toolchain lacks,
% do not invent macros -- use a minimal working preamble"), this file declares, verbatim in
% spirit and numbering behaviour, exactly the theorem-like environments OLP defines in
% sty/open-logic-book-envs.sty (thm/lem/prop/cor/defn/ex/prob/rem/note/case/axiom/conv),
% plus the shorthand accents OLP uses (`!!a`, `!!b`, `!!{...}`).
%
% Author text is untouched; only the packaging is ours.
\documentclass[10pt,a4paper]{article}

\usepackage{amsmath}
\usepackage{amssymb}
\usepackage{amsthm}
\usepackage{graphicx}
\usepackage{hyperref}
\usepackage{enumitem}
\usepackage[framemethod=default]{mdframed}

% ---- OLP shorthand accents (from sty/open-logic.sty) ----
% NB: \! is taken by LaTeX (negative thin space), so we only define the !!a / !!b families
% and the \olup helper used by them.
\newcommand*{\olup}[1]{\MakeUppercase{#1}}
\def\!!!a#1{\olup{#1}}
\def\!a#1{\olup{#1}}
\def\!b#1{\olup{#1}}

% ---- theorem-like environments, same family numbering as OLP ----
\theoremstyle{plain}
\newtheorem{thm}{Theorem}[section]
\newtheorem{lem}[thm]{Lemma}
\newtheorem{prop}[thm]{Proposition}
\newtheorem{cor}[thm]{Corollary}
\newtheorem{axiom}{Axiom}

\theoremstyle{definition}
\newtheorem{defn}[thm]{Definition}
\newtheorem{ex}[thm]{Example}
\newtheorem{prob}{Problem}[section]
\newtheorem{probd}{Problem}
\newtheorem{conv}[thm]{Convention}

\theoremstyle{remark}
\newtheorem{rem}[thm]{Remark}
\newtheorem{note}[thm]{Note}
\newtheorem{case}[thm]{Case}

% ---- environments OLP defines as wrappers ----
\newenvironment{probtag}[1]{\begin{prob}}{\end{prob}}
\newenvironment{explain}{}{}
\newenvironment{explain*}{}{}
\newenvironment{editorial}{}{}
\newenvironment{digress}{}{}
\newenvironment{digress*}{}{}
\newenvironment{history}{}{}
\newenvironment{reference}{}{}
\newcommand*{\olfileid}[3]{}
\newcommand*{\olsection}[1]{\section*{#1}}
\newcommand*{\olchapter}[3]{}
\newcommand*{\olpart}[2]{}
\newcommand*{\olimport}[1]{}
\newcommand*{\olinclude}[1]{}
\newcommand*{\olpath}[1]{}
\newcommand*{\ollangid}[1]{}
\newcommand*{\problemtag}[1]{}
\newcommand*{\problemauthor}[1]{}
\newcommand*{\ollabel}[1]{\label{#1}}
% \olref takes 0-3 OPTIONAL arguments plus one mandatory label:
%     \olref{lab}  |  \olref[from]{lab}  |  \olref[from][to]{lab}  |  \olref[a][b][c]{lab}
% A single-argument definition silently swallows "[sth]" and emits a broken reference
% (measured on the host: the text came out as "??iso]x"). The chain below consumes the
% brackets and cross-references only the mandatory argument.
\makeatletter
\newcommand{\olref}{\@ifnextchar[{\olref@i}{\olref@final}}
\def\olref@i[#1]{\@ifnextchar[{\olref@ii}{\olref@final}}
\def\olref@ii[#1]{\@ifnextchar[{\olref@iii}{\olref@final}}
\def\olref@iii[#1]{\olref@final}
\def\olref@final#1{\ref{#1}}
\makeatother
\newcommand*{\olnewpage}{\clearpage}

% ===================== OLP compatibility layer (1) =====================

\makeatletter

% ---------------- proof trees (OLP's prooftree is a bussproofs derivative) ----------------
\providecommand{\AxiomC}[1]{}
\providecommand{\UnaryInfC}[1]{}
\providecommand{\BinaryInfC}[1]{}
\providecommand{\TrinaryInfC}[1]{}
\providecommand{\QuaternaryInfC}[1]{}
\providecommand{\QuinaryInfC}[1]{}
\providecommand{\RightLabel}[1]{}
\providecommand{\LeftLabel}[1]{}
\providecommand{\doubleLine}{}
\providecommand{\noLine}{}
\providecommand{\DisplayProof}{}
\providecommand{\fCenter}{}
\newenvironment{prooftree}{}{}
\newenvironment{ltprooftree}{}{}
\newenvironment{oltableau}{}{}

% ---------------- logical notation ----------------
\providecommand{\lif}{\mathbin{\rightarrow}}
\providecommand{\lforall}{\mathord{\forall}}
\providecommand{\lexists}{\mathord{\exists}}
\providecommand{\ltrue}{\top}
\providecommand{\lfalse}{\bot}
\providecommand{\Ltrue}{\top}
\providecommand{\Lfalse}{\bot}
\providecommand{\Sequent}{\mathrel{\Rightarrow}}
\providecommand{\Proves}{\mathrel{\vdash}}
\providecommand{\Entails}{\mathrel{\models}}
\providecommand{\Subst}[3]{#1[#2/#3]}
\providecommand{\sFmla}[2]{#1(#2)}
\providecommand{\pFmla}[2]{#1(#2)}
\providecommand{\formula}{\relax}
\providecommand{\True}{\top}
\providecommand{\False}{\bot}
\providecommand{\Args}[1]{#1}
\providecommand{\Ts}{\ensuremath{\top}}
\providecommand{\Fi}{\ensuremath{\bot}}
\providecommand{\lfalsum}{\bot}

% ---------------- objects, structures, metatheory ----------------
\providecommand{\Obj}{}
\providecommand{\Struct}{\ensuremath{\mathfrak{A}}}
\providecommand{\Domain}{\ensuremath{|\mathfrak{A}|}}
\providecommand{\Assign}{\ensuremath{s}}
\providecommand{\Sat}[2]{\ensuremath{#1 \models #2}}
\providecommand{\mSat}[3]{\ensuremath{#1 \models #2[#3]}}
\providecommand{\mModel}[2]{\ensuremath{#1 \models #2}}
\providecommand{\Lang}{\ensuremath{\mathcal{L}}}
\providecommand{\Th}{\ensuremath{\mathrm{Th}}}
\providecommand{\Log}{\ensuremath{\mathrm{L}}}
\providecommand{\tuple}[1]{\langle #1 \rangle}
\providecommand{\num}[1]{\underline{#1}}
\providecommand{\cfind}[1]{\varphi_{#1}}
\providecommand{\Nat}{\ensuremath{\mathbb{N}}}
\providecommand{\Mult}{\ensuremath{\mathrm{Mult}}}
\providecommand{\Add}{\ensuremath{\mathrm{Add}}}
\providecommand{\Exp}{\ensuremath{\mathrm{Exp}}}
\providecommand{\Zero}{\ensuremath{0}}
\providecommand{\Succ}{\ensuremath{\mathrm{S}}}
\providecommand{\lambd}{\ensuremath{\lambda}}
\providecommand{\bredpar}{\mathrel{\rightarrow_\beta}}
\providecommand{\bcd}[1]{\ensuremath{\ulcorner #1 \urcorner}}
\providecommand{\Set}[1]{\ensuremath{\{#1\}}}
\providecommand{\indcase}[1]{\par\noindent\emph{Case #1}.\quad}
\providecommand{\indfrm}[1]{\par\noindent\emph{Form #1}.\quad}
\providecommand{\iftag}[1]{#1}
\providecommand{\tagitem}[1]{\item[#1]}
\providecommand{\Elim}[1]{}
\providecommand{\Intro}[1]{}
\providecommand{\LeftR}[1]{}
\providecommand{\RightR}[1]{}
\providecommand{\TAss}[1]{}
\providecommand{\TRule}[1]{}

\makeatother
\providecommand{\olpath}{.}

% ===================== OLP compatibility layer (2) =====================
% --- second compatibility block: notation and shorthands used by the frozen OLP units ---
\providecommand{\CutCS}{\mathrm{Cut}}
\providecommand{\LeftSubproofLabel}{\relax}
\providecommand{\RightSubproofLabel}{\relax}
\providecommand{\maeh}{\relax}
\providecommand{\cheight}{\relax}
\providecommand{\ande}{\mathrm{and}_E}
\providecommand{\andi}{\mathrm{and}_I}
\providecommand{\ore}{\mathrm{or}_E}
\providecommand{\ori}{\mathrm{or}_I}
\providecommand{\typeof}{\mathord{:}}
\providecommand{\pValue}[1]{#1}
\providecommand{\Sent}{\mathrm{Sent}}
\providecommand{\cardplus}{\mathbin{+}}
\providecommand{\cardtimes}{\mathbin{\cdot}}
\providecommand{\cardexpo}{\mathbin{\hat{\ }}}
\providecommand{\cardnless}{\not<}
\providecommand{\cardle}{\leq}
\providecommand{\funfromto}[3]{#1\colon #2 \to #3}
\providecommand{\funimage}[1]{\operatorname{im}(#1)}
\providecommand{\fdefined}{\mathord{\downarrow}}
\providecommand{\canonord}{\operatorname{ord}}
\providecommand{\ordtimes}{\mathbin{\cdot}}
\providecommand{\ordeq}[2]{#1 = #2}
\providecommand{\ordneq}[2]{#1 \neq #2}
\providecommand{\ordtype}[1]{\operatorname{ord}(#1)}
\providecommand{\ordsucc}{\mathord{+1}}
\providecommand{\ordunion}{\cup}
\providecommand{\onesphere}{S^1}
\providecommand{\rotationsgroup}{\mathrm{SO}(2)}
\providecommand{\equivrep}[1]{[#1]}
\providecommand{\rlexless}{\prec}
\providecommand{\ZFC}{\mathrm{ZFC}}
\providecommand{\ZF}{\mathrm{ZF}}
\providecommand{\Rat}{\mathbb{Q}}
\providecommand{\Real}{\mathbb{R}}
\providecommand{\PosInt}{\mathbb{Z}^{+}}
\providecommand{\Id}{\mathrm{id}}
\providecommand{\Bin}{\mathrm{Bin}}
\providecommand{\pto}{\rightharpoonup}
\providecommand{\dom}{\operatorname{dom}}
\providecommand{\ran}{\operatorname{ran}}
\providecommand{\len}{\operatorname{len}}
\providecommand{\concat}{\mathbin{{}^\frown}}
\providecommand{\card}[1]{\lvert #1 \rvert}
\providecommand{\Pow}{\mathcal{P}}
\providecommand{\Setabs}[2]{\{\, #1 \mid #2 \,\}}
\providecommand{\comp}{\circ}
\providecommand{\Char}{\chi}
\providecommand{\Prf}{\operatorname{Prf}}
\providecommand{\Prov}{\operatorname{Prov}}
\providecommand{\Proj}{\operatorname{proj}}
\providecommand{\Gn}{\operatorname{gn}}
\providecommand{\gn}{\operatorname{gn}}
\providecommand{\Value}[1]{\operatorname{val}(#1)}
\providecommand{\Deduce}{\mathrm{D}}
\providecommand{\Dual}{\mathrm{d}}
\providecommand{\Ax}{\mathrm{Ax}}
\providecommand{\pSat}[3]{#1 \models #2\,[#3]}
\providecommand{\pAssign}{\sigma}
\providecommand{\bforall}{\forall}
\providecommand{\bexists}{\exists}
\providecommand{\liff}{\mathbin{\leftrightarrow}}
\providecommand{\FV}{\operatorname{FV}}
\providecommand{\UnaryInf}{\ensuremath{\mid}}
\providecommand{\BinaryInf}{\ensuremath{\mid}}
\providecommand{\depth}{\operatorname{depth}}
\providecommand{\oliflabeldef}[3]{#1}
\providecommand{\TMblank}{\sqcup}
\providecommand{\TMstroke}{\mid}
\providecommand{\TMstay}{\mathrm{S}}
\providecommand{\TMright}{\mathrm{R}}
\providecommand{\TMleft}{\mathrm{L}}
\providecommand{\TMendtape}{\triangleright}
\providecommand{\umin}{\mu}
\providecommand{\Sum}{\sum}
\providecommand{\citep}[1]{\cite{#1}}
\providecommand{\citet}[1]{\cite{#1}}
\providecommand{\citeauthor}[1]{\cite{#1}}
\providecommand{\citeyear}[1]{\cite{#1}}

% ===================== OLP compatibility layer (3) =====================
% --- third compatibility block: remaining notation seen in the frozen OLP units ---
\providecommand{\iso}{\cong}
\providecommand{\nszero}{\mathrm{0}}
\providecommand{\nssucc}{\mathrm{S}}
\providecommand{\nsplus}{\mathbin{+}}
\providecommand{\nssless}{<}
\providecommand{\nsless}{<}
\providecommand{\PL}{\mathrm{PL}}
\providecommand{\Nec}{\mathord{\Box}}
\providecommand{\RK}{\mathrm{RK}}
\providecommand{\ST}{\mathrm{ST}}
\providecommand{\SSubst}[3]{#1[#2/#3]}
\providecommand{\subst}[3]{#1[#2/#3]}
\providecommand{\tagfalse}{\mathrm{F}}
\providecommand{\cutr}{\mathrm{Cut}}
\providecommand{\cardneq}{\neq}
\providecommand{\issucc}[1]{\operatorname{succ}(#1)}
\providecommand{\tup}[1]{\langle #1 \rangle}
\providecommand{\Var}{\mathrm{Var}}
\providecommand{\Bew}{\operatorname{Bew}}
\providecommand{\Con}{\operatorname{Con}}
\providecommand{\Cons}{\operatorname{Cons}}
\providecommand{\pDeduce}{\mathrm{D}}
\providecommand{\psat}{\models}
\providecommand{\den}[1]{\llbracket #1 \rrbracket}
\providecommand{\Sem}[2]{\llbracket #1 \rrbracket_{#2}}
\providecommand{\DomainOf}[1]{\lvert #1 \rvert}
\providecommand{\Form}{\mathrm{Form}}
\providecommand{\Terms}{\mathrm{Term}}
\providecommand{\Atom}{\mathrm{Atom}}
\providecommand{\FVars}[1]{\operatorname{FV}(#1)}
\providecommand{\fv}[1]{\operatorname{fv}(#1)}
\providecommand{\tuplep}[1]{\langle #1 \rangle}
\providecommand{\restrict}[2]{#1 \upharpoonright #2}
\providecommand{\sto}{\rightarrow}
\providecommand{\FmlaSet}[1]{\mathrm{Fmla}_{#1}}
\providecommand{\Thm}{\mathrm{Thm}}
\providecommand{\Eval}{\operatorname{val}}
\providecommand{\blank}{\sqcup}
\providecommand{\blankTape}{\sqcup}
\providecommand{\Consis}{\operatorname{Con}}
\providecommand{\drs}{\mathrm{d}}
\providecommand{\VRule}{\mid}
\providecommand{\Bet}{\mathrm{B}}
\providecommand{\Seg}[1]{\mathrm{Seg}_{#1}}
\providecommand{\Enums}{\mathrm{En}}
\providecommand{\Z}{\mathbb{Z}}
\providecommand{\N}{\mathbb{N}}
\providecommand{\abs}[1]{\lvert #1 \rvert}
\providecommand{\floor}[1]{\lfloor #1 \rfloor}
\providecommand{\ceil}[1]{\lceil #1 \rceil}

% ===================== OLP compatibility layer (4) =====================
% --- fourth compatibility block: harvested from the per-candidate compile logs ---
\providecommand{\beredpar}{\mathrel{\rightarrow_\beta}}
\makeatletter
\def\bered{\@ifnextchar[{\bered@i}{\mathrel{\rightarrow_\beta}}}
\def\bered@i[#1]{\mathrel{\rightarrow_\beta}}
\def\bred{\@ifnextchar[{\bred@i}{\mathrel{\rightarrow_\beta}}}
\def\bred@i[#1]{\mathrel{\rightarrow_\beta}}
\makeatother
\providecommand{\beredone}{\mathrel{=_\beta}}
\providecommand{\bredone}{\mathrel{=_\beta}}
\providecommand{\xredone}{\mathrel{=_\beta}}
\providecommand{\aconvone}{\mathrel{=_\alpha}}
\providecommand{\cardless}{<}
\providecommand{\tagrefs}[1]{#1}
\providecommand{\ordplus}{\mathbin{+}}
\providecommand{\funrestrictionto}[2]{#1 \upharpoonright #2}
\providecommand{\fundefined}{\mathord{\uparrow}}
\providecommand{\Trm}{\mathrm{Trm}}
\providecommand{\OProv}{\operatorname{Prov}}
\providecommand{\OPrf}{\operatorname{Prf}}
\providecommand{\aeq}{\mathrel{=}}
\providecommand{\mClass}[1]{[#1]}
\providecommand{\equivclass}[1]{[#1]}
\providecommand{\pheight}{\operatorname{ht}}
\providecommand{\phy}{\operatorname{ht}}
\providecommand{\Cut}{\mathrm{Cut}}
\providecommand{\QuantRank}{\operatorname{qr}}
\providecommand{\elemequiv}{\mathrel{\equiv}}
\providecommand{\PIso}{\mathrm{PIso}}
\providecommand{\Theory}{\mathrm{Th}}
\providecommand{\supstrict}{\subsetneq}
\providecommand{\trcl}{\operatorname{trcl}}
\providecommand{\cref}[1]{\ref{#1}}
\providecommand{\Discharge}{\mathrm{Dis}}
\providecommand{\Weakening}{\mathrm{Wk}}
\providecommand{\ZFminus}{\mathrm{ZF}^{-}}
\makeatletter
\renewcommand{\Dual}{\@ifnextchar[{\Dual@i}{\ifmmode\mathrm{d}\else$\mathrm{d}$\fi}}
\def\Dual@i[#1]{\ifmmode\mathrm{d}[#1]\else$\mathrm{d}[#1]$\fi}
\makeatother
\providecommand{\RightR}{\mathrm{R}}
\providecommand{\LeftR}{\mathrm{L}}
\providecommand{\Intro}{\mathrm{I}}
\providecommand{\Axiom}{\mathrm{Ax}}
\providecommand{\formula}[1]{#1}
\providecommand{\stared}{\ast}
\providecommand{\oset}[2]{\{\, #1 \mid #2 \,\}}
\providecommand{\setof}[2]{\{\, #1 \mid #2 \,\}}
\providecommand{\pow}[1]{\mathcal{P}(#1)}
\providecommand{\ordpair}[2]{\langle #1, #2 \rangle}
\providecommand{\ordtriple}[3]{\langle #1, #2, #3 \rangle}
\providecommand{\union}{\cup}
\providecommand{\inter}{\cap}
\providecommand{\seq}[1]{\langle #1 \rangle}
\providecommand{\SuchThat}{\mid}
\providecommand{\where}{\mid}
\providecommand{\TVar}{\mathrm{TVar}}
\providecommand{\Bool}{\mathrm{Bool}}
\providecommand{\Val}{\operatorname{Val}}
\providecommand{\TL}[1]{\mathrm{TL}}
\providecommand{\xstared}{\ast}

% ===================== OLP compatibility layer (5) =====================
% --- fifth compatibility block: rescues the quarantined OLP candidates ---
\providecommand{\ident}{\operatorname{id}}
\providecommand{\setrank}{\operatorname{rank}}
\providecommand{\cardfont}{\mathfrak{c}}
\providecommand{\Complement}{\overline}
\providecommand{\QR}{\mathrm{QR}}
\providecommand{\becd}{\ensuremath{\ulcorner}}
\providecommand{\disjointsum}{\mathbin{\oplus}}
\providecommand{\cardeq}[2]{#1 = #2}
\providecommand{\xred}{\mathrel{\rightarrow}}
% --- sixth compatibility block: names reported by real compile logs ---
\providecommand{\ORefut}{\operatorname{Refut}}
\providecommand{\ORProv}{\operatorname{Prov}}
\providecommand{\cutrank}{\operatorname{cr}}
\providecommand{\Refut}{\operatorname{Refut}}
\providecommand{\DischargeRule}{\mathrm{Dis}}
\providecommand{\Contraction}{\mathrm{C}}
\providecommand{\bottomAlignProof}{\relax}
\providecommand{\FalseCl}{\bot}
\providecommand{\red}{\mathrel{\rightarrow}}
\providecommand{\varAssign}{\sigma}
\providecommand{\defiff}{\mathrel{\leftrightarrow}}
\providecommand{\FalseInt}{\bot}
\providecommand{\fn}{\operatorname{fn}}
\providecommand{\dotted}{\cdot}

% --- tag-macro semantic fix (2026-10-13, final): \iftag CHOOSES a branch and the OLP proof body is
% always the SECOND group (\iftag{tag}{body}, {tag}{body}{c}, {tag}{body}{c}{d}); a definition that
% prints the tag literal drops the proof body. Bodies contain blank lines, hence \long.
\makeatletter
\providecommand{\iftag}{}
\long\def\iftag#1#2{\@ifnextchar\bgroup{\iftag@three{#2}}{\iftag@stop{#2}}}
\long\def\iftag@three#1#2{\iftag@stop{#1}\@ifnextchar\bgroup{\iftag@eat}{\relax}}
\long\def\iftag@eat#1{}
\long\def\iftag@stop#1{#1}
\long\def\oliflabeldef#1#2{\@ifnextchar\bgroup{\@gobble}{\iftag@stop{#2}}}
\providecommand{\tagrefs}[1]{}
\providecommand{\settag}[1]{}
\makeatother
\begin{document}

% ==== OLP file: content/proof-theory/cut-elimination/ce-largest.tex ====
% Part: proof-theory
% Chapter: cut-elimination
% Section: ce-largest



\olfileid{pt}{cut}{inv}

\olsection{Removing Largest Cuts}

The proof of \olref[seq][inv]{lem:G3c-invert} shown that if \Log{G3c}
\olup{prove}s the conclusion of a logical rule (other than
\RightR{\lexists} and \LeftR{\lforall}), it also \olup{prove}s each of
the premises of the rule without increasing the \olup{height} of the
proof. We can do better and show that this holds also for $\Log{G3c} +
\CutCS$.

As before, we take the \emph{cut rank} of a \CutCS{} inference to be
the depth of its cut \olup{formula}.

\begin{lem}\ollabel{lem:inv-G3c-cut}
If $\Log{G3c} + \CutCS$ \olup{prove}s the conclusion of a logical
rule~$R$ other than \RightR{\lexists} or \LeftR{\lforall} using \olup{proof}~$\pi$ then it also \olup{prove}s each premise
with \olup{proof}~$\pi'$ (or \olup{proof}s~$\pi'$, $\pi''$ depending on
whether the rule has one or two premises). Moreover,
(a)~$\pheight{\pi'}\le\pheight{\pi}$ (and
$\pheight{\pi''}\le\pheight{\pi}$) and (b)~the number of \CutCS{}
inferences as well as their ranks is the same in~$\pi'$ (and $\pi''$)
as in~$\pi$.
\end{lem}

\begin{proof}
By inspection of the proofs of
\cref{pt:seq:inv:lem:G3c-invert,pt:seq:inv:lem:invert-quant}. That
proof was by induction on~$\pheight{\pi}$. In the base case, we
replaced an axiom by another axiom. So conditions (a) and~(b) are
trivially satisfied. If the last rule of~$\pi$ was~$R$, its immediate
sub-\olup{proof}s are the \olup{proof}s~$\pi_i$ we want, and (a) and~(b) are
both satisfied. In all other cases, we applied the inductive
hypothesis to the immediate sub-\olup{proof}s of~$\pi$, i.e., those
leading to the premise(s) of the last inference~$R$, and then applied
that same rule $R$ to the results. The conditions (a) and~(b) hold for
the immediate sub-\olup{proof}s of the new proof, and so (a) and~(b) hold
for the entire proof. It remains to verify the case where the last
inference is~\CutCS. We do this again for one example only: suppose
$\pi$ is \olup{proof} of the conclusion $B \land C, \Gamma \Sequent
\Delta$ of \LeftR{\land}, and ends in~\CutCS:
\[
\AxiomC{}
\RightLabel{$\pi_1$}
\Deduce$B \land C, \Gamma \fCenter \Delta, A$
\AxiomC{}
\RightLabel{$\pi_2$}
\Deduce$A, B \land C, \Gamma \fCenter \Delta$
\RightLabel{\CutCS}
\BinaryInf$B \land C, \Gamma \fCenter \Delta$
\DisplayProof
\]
The induction hypothesis applies to $\pi_1$ and~$\pi_2$ (which are
also \olup{proof}s of instances of the conclusion of $\LeftR{\land}$) and
yields \olup{proof}s $\pi_1'$ and~$\pi_2'$ of $B, C, \Gamma \fCenter
\Delta, A$ and $A, B, C, \Gamma \fCenter \Delta$,
respectively. We combine these using \CutCS{} to obtain~$\pi'$:
\[
\AxiomC{}
\RightLabel{$\pi_1'$}
\Deduce$B, C, \Gamma \fCenter \Delta, A$
\AxiomC{}
\RightLabel{$\pi_2'$}
\Deduce$A, B, C, \Gamma \fCenter \Delta$
\RightLabel{\CutCS}
\BinaryInf$B, C, \Gamma \fCenter \Delta$
\DisplayProof
\]
Clearly, (a) and~(b) are satisfied.
\end{proof}

Note that this would not work if we used \Cut{} instead of \CutCS, as
then the conclusion of the last \olup{proof} would have two copies of
$B, C, \Gamma$ in the antecedent and two copies of $\Delta$ in the
succedent. We would have to appeal to the admissibility of
contraction, but the proof of \olref[seq][inv]{prop:G3c-cont-adm} uses
the invertibility lemma. We'd thus be arguing in a circle.

We can use \cref{pt:cut:inv:lem:inv-G3c-cut} to give an alternative
proof of cut elimination. In this proof we do not successively remove
topmost occurrences of \CutCS{} inferences, but occurrences of \CutCS{}
inferences of maximal rank. That is, we start with \olup{proof}~$\pi$ in
$\Log{G3c} + \CutCS$ and remove a cut anywhere in it (possibly
replacing it with cuts of smaller rank)---and then keep doing that
until no cuts remain. The only condition is that above the cut we deal
with, no cuts of the same or higher rank occur.

\begin{lem}\ollabel{lem:max-cut-red-G3c} If $\pi$ is \olup{proof} in
$\Log{G3c}+\CutCS$ ending in a~\CutCS{} inference of rank~$n$ and
otherwise using only rules of~$\Log{G3c}$ and \CutCS{} inferences of
rank $<n$, then there is a proof~$\pi'$ of the same end-sequent
in~$\Log{G3c} + \CutCS$ where every \CutCS{} inference has rank~$<n$.
\end{lem}

\begin{proof}
The \olup{proof}~$\pi$ ends in:
\[
\AxiomC{}
\RightLabel{$\pi_1$}
\Deduce$\Gamma \fCenter \Delta, A$
\AxiomC{}
\RightLabel{$\pi_2$}
\Deduce$A, \Gamma \fCenter \Delta$
\RightLabel{\CutCS}
\BinaryInf$\Gamma \fCenter \Delta$
\DisplayProof
\]
We distinguish cases according to the form of~$A$.

Suppose $A$ is atomic. Observe that by the condition that all
\CutCS{} inferences in~$\pi_1$ and $\pi_2$ must have rank $<
\depth{A} = 0$, there can be no \CutCS{} inferences in~$\pi_1$
or~$\pi_2$. We show that $\Gamma \Sequent \Delta$ has a cut-free proof
by induction on $\pheight{\pi_2}$. If $\pheight{\pi_2} = 0$, then $A,
\Gamma \Sequent \Delta$ is an axiom. If $A$ is not principal, some
atomic $C$ is \olup{element} of both~$\Gamma$ and $\Delta$, and $\Gamma
\Sequent \Delta$ itself is an axiom. If $A$ is principal, then
$\Delta = \Delta, A$. In this case, $\pi_1$ ends in $\Gamma \Sequent
\Delta', A, A$. We obtain a cut-free proof of $\Gamma \Sequent
\Delta', A$ by applying \olref[seq][inv]{prop:G3c-cont-adm}
(admissibility of contraction). If $\pheight{\pi_2} > 0$ it ends in a
rule~$R$. Take a premise of that rule: it has the form $A, \Gamma',
\Pi \Sequent \Delta', \Lambda$, where $\Pi$ and $\Lambda$ are the
active \olup{formula}s of~$R$. \olref[seq][adm]{prop:weak-G3c-adm}
applied to~$\pi_1$ and $\pi_2$ yields \olup{proof}s of $\Gamma, \Gamma',
\Pi \Sequent \Delta, \Delta', \Lambda, A$ and $A, \Gamma, \Gamma',
\Pi \Sequent \Delta, \Delta', \Lambda$, to which the induction
hypothesis applies. This gives us \olup{proof} of $\Gamma, \Gamma', \Pi
\Sequent \Delta, \Delta', \Lambda$, for each premise of~$R$. Applying
$R$ yields $\Gamma, \Gamma \Sequent \Delta, \Delta$, and
\olref[seq][inv]{prop:G3c-cont-adm} provides the cut-free proof of
$\Gamma \Sequent \Delta$.

If $A$ is not atomic, we distinguish cases according to the \olup{main
operator} of~$A$. In each case, we apply \olref{lem:inv-G3c-cut} to
obtain \olup{proof}s of the premises of the left- and right-rules with
$A$ as principal formula. We then proceed as in case~E of the proof
of \olref[top]{lem:cut-adm-G3c}.
\begin{enumerate}
  \item $A \ident \lnot B$. The \olup{proof}s $\pi_1$ and $\pi_2$ end
  in $\Gamma \Sequent \Delta, \lnot B$ and $\lnot B, \Gamma \Sequent
  \Delta$. By \olref{lem:inv-G3c-cut}, there are \olup{proof}s $\pi_1'$
  and $\pi_2'$ of $B, \Gamma \Sequent
  \Delta$ and $\Gamma \Sequent \Delta, B$. The \olup{proof}~$\pi'$ is:
  \[
  \AxiomC{}
  \RightLabel{$\pi_2'$}
  \Deduce$\Gamma \fCenter \Delta, B$
  \AxiomC{}
  \RightLabel{$\pi_1'$}
  \Deduce$B, \Gamma \fCenter \Delta$
  \RightLabel{\CutCS}
  \BinaryInf$\Gamma \fCenter \Delta$
  \DisplayProof
  \]
  \item $A \ident (B \lor C)$. The \olup{proof}s $\pi_1$ and $\pi_2$
  end in $\Gamma \Sequent \Delta, B \lor C$ and $B \lor C, \Gamma
  \Sequent \Delta$. By \olref{lem:inv-G3c-cut} (inversion for for
  \RightR{\lor} applied to~$\pi_1$), there is \olup{proof} $\pi_1'$ of
  $\Gamma \Sequent \Delta, B, C$. By \olref{lem:inv-G3c-cut}
  (inversion for \LeftR{\lor} applied to~$\pi_2$), there is \olup{proof}
  $\pi_2'$ of $B, \Gamma \Sequent \Delta$ and \olup{proof} $\pi_2''$ of
  $C, \Gamma \Sequent \Delta$. By applying
  \olref[seq][adm]{prop:weak-G3c-adm} to $\pi_2''$ we also get
  \olup{proof}~$\pi_2'''$ of $C, \Gamma \Sequent \Delta, B$. The
  \olup{proof}~$\pi'$ is:
  \[
  \AxiomC{}
  \RightLabel{$\pi_1'$}
  \Deduce$\Gamma \fCenter \Delta, B, C$
  \AxiomC{}
  \RightLabel{$\pi_2'''$}
  \Deduce$C, \Gamma \fCenter \Delta, B$
  \RightLabel{\CutCS}
  \BinaryInf$\Gamma \fCenter \Delta, B$
  \AxiomC{}
  \RightLabel{$\pi_2'$}
  \Deduce$B, \Gamma \fCenter \Delta$
  \RightLabel{\CutCS}
  \BinaryInf$\Gamma \fCenter \Delta$
  \DisplayProof
  \]
  \item $A \ident (B \land C)$. Exercise.
  \item $A \ident (B \lif C)$. The \olup{proof}s $\pi_1$ and $\pi_2$
  end in $\Gamma \Sequent \Delta, B \lif C$ and $B \lif C, \Gamma
  \Sequent \Delta$. By \olref{lem:inv-G3c-cut} (inversion for
  \RightR{\lif} applied to~$\pi_1$), there is \olup{proof} $\pi_1'$ of
  $B, \Gamma \Sequent \Delta, C$. By \olref{lem:inv-G3c-cut}
  (inversion for \LeftR{\lif} applied to~$\pi_2$), there is
  \olup{proof} $\pi_2'$ of $\Gamma \Sequent \Delta, B$ and \olup{proof}
  $\pi_2''$ of $C, \Gamma \Sequent \Delta$. By applying
  \olref[seq][adm]{prop:weak-G3c-adm} to $\pi_2''$ we also get
  \olup{proof}~$\pi_2'''$ of $C, B, \Gamma \Sequent \Delta$. The
  \olup{proof}~$\pi'$ is:
  \[
  \AxiomC{}
  \RightLabel{$\pi_2'$}
  \Deduce$\Gamma \fCenter \Delta, B$
  \AxiomC{}
  \RightLabel{$\pi_1'$}
  \Deduce$B, \Gamma \fCenter \Delta, C$
  \AxiomC{}
  \RightLabel{$\pi_2'''$}
  \Deduce$C, B, \Gamma \fCenter \Delta$
  \RightLabel{\CutCS}
  \BinaryInf$B, \Gamma \fCenter \Delta$
  \RightLabel{\CutCS}
  \BinaryInf$\Gamma \fCenter \Delta$
  \DisplayProof
  \]
  \item $A \ident \lforall[x][B(x)]$. The \olup{proof}s $\pi_1$ and
  $\pi_2$ end in $\Gamma \Sequent \Delta, \lforall[x][B(x)]$ and
  $\lforall[x][B(x)], \Gamma \Sequent \Delta$. By
  \olref{lem:inv-G3c-cut}, there is \olup{proof}~$\pi_1'$ of $\Gamma
  \Sequent \Delta, B(c)$.

  Now consider the \olup{proof}~$\pi_2$. It ends in $\lforall[x][B(x)],
  \Gamma \Sequent \Delta$. Moving upward from the end-sequent
  of~$\pi_2$, mark every occurrence of $\lforall[x][B(x)]$ on the
  left of a sequent that ``leads to'' the occurrence of
  $\lforall[x][B(x)]$ in the end-sequent of~$\pi_2$, that is:
  (a)~Mark $\lforall[x][B(x)]$ in the end-sequent of~$\pi_2$. (b)~If
  $\lforall[x][B(x)]$ occurs in the context of the conclusion of a
  rule and is marked, mark a corresponding occurrence in the
  premise(s) of the rule. (c)~If $\lforall[x][B(x)]$ is principal in
  the conclusion of a \LeftR{\lforall} rule and is marked, mark the
  corresponding active occurrences of $\lforall[x][B(x)]$ and $B(t)$
  in the premise.
  
  Now consider all these marked occurrences of $\lforall[x][B(x)]$.
  None of them are active in a rule other than \LeftR{\lforall} (only
  active \olup{formula}s of \LeftR{\lforall} get marked). Delete all
  marked occurrences of $\lforall[x][B(x)]$ in $\pi_2$.
  If this is a correct \olup{proof}, the marked occurrences of
  $\lforall[x][B(x)]$ were never active; they all came directly from
  axioms. The \olup{proof} ends in~$\Gamma \Sequent \Delta$ and we can
  replace $\pi$ by it. We have a removed a cut of maximal rank.

  Otherwise, what we have is not a correct proof because we have
  deleted \olup{formula}s $\lforall[x][B(x)]$ which were active in some
  \LeftR{\lforall} inferences. We have turned these into
  ``inferences'' of the form
  \[
  \AxiomC{}
  \RightLabel{$\pi_t$}
  \Deduce$B(t), \Pi \fCenter \Lambda$
  \RightLabel{\LeftR{\lforall}}
  \UnaryInf$\Pi \fCenter \Lambda$
  \DisplayProof
  \]
  Replace each topmost such ``inference'' by
  \[
  \AxiomC{}
  \RightLabel{$\Subst{\pi_1'}{t}{c}[\Pi \Sequent \Lambda]$}
  \Deduce$\Gamma, \Pi \fCenter \Delta, \Lambda, B(t)$
  \AxiomC{}
  \RightLabel{$\pi_t[\Gamma \Sequent \Delta]$}
  \Deduce$B(t), \Gamma, \Pi \fCenter \Delta, \Lambda$
  \RightLabel{\CutCS}
  \BinaryInf$\Gamma, \Pi \fCenter \Delta, \Lambda$
  \DisplayProof
  \]
  and add $\Gamma$ to the left and $\Delta$ to the right of every
  sequent below it. Repeat until all \LeftR{\lforall} ``inferences''
  with a marked $B(t)$ in the premise have been replaced by a
  \CutCS{} inference on some~$B(t)$. The end-sequent of the resulting
  \olup{proof} (and it now is a correct \olup{proof}) is of the form
  $\Gamma, \dots, \Gamma \Sequent \Delta, \dots, \Delta$. Apply
  admissibility of contraction to turn it into \olup{proof} of~$\Gamma
  \Sequent \Delta$ and replace $\pi$ with it. We have replaced one cut
  on $\lforall[x][B(x)]$ with (possibly many) cuts on $B(t)$, but
  these are all of lower rank. Any cuts already occurring in $\pi_1$ or
  $\pi_2$ remain present, but they are also all of lower rank.
\end{enumerate}
\end{proof}

\begin{prob}
Verify the case where $A \ident (B \land C)$ in the proof of
\olref{lem:inv-G3c-cut}. That is, show that if there is
\olup{proof}~$\pi$ ending in 
\[
\AxiomC{}
\RightLabel{$\pi_1$}
\Deduce$\Gamma \fCenter \Delta, B \land C$
\AxiomC{}
\RightLabel{$\pi_2$}
\Deduce$B \land C, \Gamma \fCenter \Delta$
\RightLabel{\CutCS}
\BinaryInf$\Gamma \fCenter \Delta$
\DisplayProof
\]
with $\pi_1$ and $\pi_2$ containing only cuts of rank $<\depth{B
\land C}$ then there is \olup{proof}~$\pi'$ of $\Gamma \Sequent \Delta$
containing only cuts of rank $<\depth{B \land C}$.
\end{prob}

% \begin{prob}
% Verify the case where $A \ident \lexists[x][B(x)]$ in the proof of
% \olref{lem:inv-G3c-cut}.
% \end{prob}


\olref[top]{lem:cut-adm-G3c} establishes that we can replace a \CutCS{}
inference of highest rank in \olup{proof} by \CutCS{} inferences of
lower rank. We can again use this lemma to show that we can eliminate
any number of \CutCS{} inferences from \olup{proof} by successively
applying it to the topmost sub-\olup{proof}s ending in maximal~\CutCS.

\begin{thm}\ollabel{thm:cut-elim-G3c-largest}
Any \olup{proof} $\pi$ in $\Log{G3c} + \CutCS$ can be transformed into a proof in~\Log{G3c}.
\end{thm}

\begin{proof}
  First we prove that all cuts of maximal rank in~$\pi$ can be
  removed. Let $d$ be the maximal rank of cuts occuring in~$\pi$, and
  let $n$ be the number of cuts of rank~$d$. The proof is by induction
  on~$n$. If $n=0$ there is nothing to prove. If $n>0$, pick a topmost
  cut of rank~$d$, and let $\pi_1$ be the sub-\olup{proof} ending in it.
  By \olref{lem:max-cut-red-G3c}, there is \olup{proof}~$\pi_1'$ of the
  same end-sequent with all cuts of rank $<d$. Replace $\pi_1$ by
  $\pi_1'$ in~$\pi$. This results in \olup{proof} containing $n-1$ cuts
  of rank~$d$, and the induction hypothesis applies.

  Now the result follows by induction on~$d$. If $d=0$, all cuts are
  atomic, and by the previous result can all be removed. If $d>0$,
  the previous result yields a proof in which all cuts of rank $d$
  have been replaced by cuts of rank~$<d$. Since $d$ was the maximum
  rank of cuts in~$\pi$, we have proof in which the maximal rank of
  cuts is $<d$, and the induction hypothesis applies.
\end{proof}

% ==== OLP file: content/proof-theory/cut-elimination/ce-topmost.tex ====
% Part: proof-theory
% Chapter: cut-elimination
% Section: ce-topmost



\olfileid{pt}{cut}{top}

\olsection{Removing Topmost Cuts}

\begin{defn}
Consider \olup{proof}~$\pi$ ending in a \CutCS{} rule,
\[
\AxiomC{}
\RightLabel{$\pi_1$}
\Deduce$\Gamma \fCenter \Delta, A$
\AxiomC{}
\RightLabel{$\pi_2$}
\Deduce$A, \Gamma \fCenter \Delta$
\RightLabel{\CutCS}
\BinaryInf$\Gamma \fCenter \Delta$
\DisplayProof
\]
and containing no other \CutCS{} inferences. The \emph{cut height}
of~$\pi$ is $\cheight{\pi} = \pheight{\pi_1} + \pheight{\pi_2}$. The
\emph{cut rank} of~$\pi$ is $\cutrank{\pi} = \depth{A}$.
\end{defn}

\begin{lem}\ollabel{lem:cut-adm-G3c}
The \CutCS{} rule is admissible in~\Log{G3c}. In other words, if $\pi$
is \olup{proof} ending in a single~\CutCS{} and otherwise using only rules
of~$\Log{G3c}$, then there is \olup{proof}~$\pi'$ of the same end-sequent in~\Log{G3c}.
\end{lem}

\begin{proof}
The proof is by double induction on cut \olup{height} and rank. The
\olup{proof}~$\pi$ ends in:
\[
\AxiomC{}
\RightLabel{$\pi_1$}
\Deduce$\Gamma \fCenter \Delta, A$
\AxiomC{}
\RightLabel{$\pi_2$}
\Deduce$A, \Gamma \fCenter \Delta$
\RightLabel{\CutCS}
\BinaryInf$\Gamma \fCenter \Delta$
\DisplayProof
\]

Induction basis: $\cheight{\pi} = 0$ and $\cutrank{\pi} = 0$. Since
$\cheight{\pi} = \pheight{\pi_1} + \pheight{\pi_2} = 0$, both $\Gamma
\Sequent \Delta, A$ and $A, \Gamma \Sequent \Delta$ are axioms.
There are two cases: (a)~either $A$ is not principal in one of them,
or (b)~$A$ is principal in both of them. In both cases we'll show
below that $\Gamma \Sequent \Delta$ must be an axiom (and we don't
need the inductive hypothesis to do that).

We'll now distinguish a number of possible cases according to whether
$\pi_1$ and $\pi_2$ are axioms or end in an inference, and whether
$A$~is principal in that inference or not. Cases A and~B cover the
induction basis. In the inductive step we assume $\cheight{\pi} > 0$
or $\cutrank{\pi} > 0$ (or both). We can use the inductive hypothesis,
namely that the last~\CutCS{} can be removed from any \olup{proof} which
ends in a single~\CutCS{} which either has cut rank $= \cutrank{\pi}$
and a cut \olup{height} $< \cheight{\pi}$, or has cut rank~$<
\cutrank{\pi}$. The case where $\cheight{\pi} = 0$ (both premises are
axioms) are covered by cases~A and~B. The cases where $\cheight{\pi} >
0$ are covered by cases C--D.

\paragraph{Case A:}
Some $B$ occurs in both $\Gamma$ and~$\Delta$. 
Then $\Gamma \Sequent \Delta$ is an axiom (if $B$~is atomic), or has
\olup{proof}~$\pi'$ in \Log{G3c} without \CutCS{} by
\olref[seq][ex]{prop:G3c-axioms}. This covers case~(a) of the
induction basis, since if $A$ is not principal in $\Gamma
\Sequent \Delta, A$ or not principal in $A, \Gamma \Sequent \Delta$
and these are axioms, then some other atomic~$B$ must occur in both
$\Gamma$ and~$\Delta$. It also covers the inductive step if one of the
premises is an axiom in which $A$ is not principal (even if the other
premise is the conclusion of a rule).

\paragraph{Case B:} Both premises are axioms, and $A$ is principal
and hence atomic. Consider the left premise, $\Gamma \Sequent
\Delta, A$. Since $A$ is principal, it must occur in $\Gamma$
(otherwise $\Gamma \Sequent \Delta, A$ wouldn't be an axiom). Similarly,
$A$ must occur in $\Delta$, or else $A, \Gamma \Sequent \Delta$
wouldn't be an axiom. Thus $A$ is atomic and occurs in both $\Gamma$
and $\Delta$, i.e., $\Gamma \Sequent \Delta$ is an axiom. This covers
case~(b) of the induction basis.

\begin{center}
Alternative cases A and B
\end{center}

\paragraph{Case A:} One of the premises of the cut is an axiom of the
form $C, \Gamma' \Sequent \Delta', C$ with $C$~atomic. If $C$ is
not the cut \olup{formula}~$A$, then $C$ must occur in both $\Gamma$
and~$\Delta$. In that case, $\Gamma \Sequent \Delta$ is an axiom and
we can take it as~$\pi'$. If $C \ident A$ is the cut \olup{formula}
then either (if the left premise is an axiom) $A \in \Gamma$ and
$\Gamma = \Gamma', A$, or (if the right premise is an axiom) $A \in
\Delta$ and $\Delta = \Delta', A$. In the first case, $\pi_2$ is
\olup{proof} of $A, A, \Gamma' \Sequent \Delta$ and by admissibility
of \LeftR{\Contraction} (\olref[seq][inv]{prop:G3c-cont-adm}) there is
\olup{proof} of $\Gamma \Sequent \Delta$. In the second case, $\pi_2$ is
\olup{proof} of $\Gamma \Sequent \Delta', A, A$ and by admissibility
of \RightR{\Contraction} we get \olup{proof} of $\Gamma \Sequent
\Delta$.

\paragraph{Case B:} One of the premises is an axiom of the
form~$\lfalse, \Gamma' \Sequent \Delta'$. If the left premise is such
an axiom, then $\lfalse \in \Gamma$ and $\Gamma \Sequent \Delta$ is
also an axiom. If the right premise is such an axiom, then either
$\lfalse \in \Gamma$ (in which case $\Gamma \Sequent \Delta$ again is
an axiom), or $A \ident \lfalse$. In the latter case, $\pi_1$ is
\olup{proof} of $\Gamma \Sequent \Delta, \lfalse$. We can turn it into
\olup{proof} of~$\Gamma \Sequent \Delta$ by deleting one copy
of~$\lfalse$ from the consequent of every sequent in~$\pi_1$.
(Exercise: verify this by induction on~$\pheight{\pi_1}$.)

Now suppose that neither case~A nor case~B applies. In the remaining
cases, $\cheight{\pi} > 0$, i.e., at least one of the premises is the
conclusion of an inference.

\paragraph{Cases C and D: Permuting cuts.}
The possibilities considered in cases~C and~D all have a common
pattern. We begin with \olup{proof} ending in~\CutCS{} where one premise
is obtained by a rule~$R$ in which the cut \olup{formula}~$A$ is not
principal (some other \olup{formula}~$D$ is). In this \olup{proof},
\CutCS{} follows the rule~$R$, and schematically looks as follows:
\[
\AxiomC{}
\RightLabel{$\pi_1$}
\Deduce$\Gamma \fCenter \Delta', D, A$
\AxiomC{}
\RightLabel{$\pi_3$}
\Deduce$A, \Gamma, \Pi \fCenter \Delta', \Lambda$
\RightLabel{$R$}
\UnaryInf$A, \Gamma \fCenter \Delta', D$
\RightSubproofLabel{$\pi_2$}
\RightLabel{\CutCS}
\BinaryInf$\Gamma \fCenter \Delta', D$
\DisplayProof
\]
This only covers the case where $R$~is a right rule with one premise,
$A$~is not the principal \olup{formula} of~$R$, and the \olup{formula}~$D$ which
\emph{is} principal occurs in the succedent of the right premise
of~\CutCS. If $D$~occurs in the antecedent (i.e., $R$ is a left rule)
or in the left premise of~\CutCS, it would look different, of course.
The \olup{formula}s in $\Pi$ and~$\Lambda$ represent the active
\olup{formula}s of rule~$R$.

We turn this order around and consider \olup{proof}
where $R$ follows~\CutCS{}:
\[
\AxiomC{}
\RightLabel{$\pi_1'$}
\Deduce$\Gamma, \Pi \fCenter \Delta', D, \Lambda, A$
\AxiomC{}
\RightLabel{$\pi_3'$}
\Deduce$A, \Gamma, \Pi \fCenter \Delta', D, \Lambda$
\RightLabel{\CutCS}
\BinaryInf$\Gamma, \Pi \fCenter \Delta', D, \Lambda$
\RightLabel{$R$}
\UnaryInf$\Gamma, \fCenter \Delta', D, D$
\DisplayProof
\]
We use \olup{height}-preserving weakening to obtain $\pi_1'$ and $\pi_3'$
from $\pi_1$ and $\pi_3$, respectively.

Because we switch the order of \CutCS{} and~$R$, this is referred to
as ``permuting a cut upwards across~$R$.'' Doing so decreases the
height of the \olup{proof} ending in the right premise of~\CutCS, and so
decreases the cut \olup{height}. That is, the sub-\olup{proof}s ending in the
new \CutCS{} can be assumed to have \olup{height} no larger than
$\pi_1$ and $\pi_3$, so the cut \olup{height} is reduced (we've removed
one inference on the right). Now the induction hypothesis applies,
which tells us that the \CutCS{} inference can be removed. Finally, we
remove the extra copy of~$D$ using the admissibility
of~$\RightR{\Contraction}$.

\paragraph{Case C:}
At least one of $\pi_1$ and $\pi_2$ ends in a one-premise rule~$R$ in
which $A$~is not principal. There are a total of 18 cases, depending
on whether $A$~is not principal in the left or the right premise
of~\CutCS, and depending on which of the nine one-premise rules
(\LeftR{\lnot}, \RightR{\lnot}, \RightR{\lor}, \LeftR{\land},
\RightR{\lif}, and the four quantifier rules) the rule~$R$ is. We only
consider two examples: where $A$~is not principal in the last
inference of~$\pi_2$, which ends in~$\RightR{\lif}$ or
in~\LeftR{\lexists}. The other cases are analogous and left as
exercises.

Suppose $\pi$ ends in:
\[
\AxiomC{}
\RightLabel{$\pi_1$}
\Deduce$\Gamma \fCenter \Delta', B \lif C, A$
\AxiomC{}
\RightLabel{$\pi_3$}
\Deduce$A, B, \Gamma \fCenter \Delta', C$
\RightLabel{$\RightR{\lif}$}
\UnaryInf$A, \Gamma \fCenter \Delta', B \lif C$
\RightSubproofLabel{$\pi_2$}
\RightLabel{\CutCS}
\BinaryInf$\Gamma \fCenter \Delta', B \lif C$
\DisplayProof
\]
where $\Delta = \Delta', B \lif C$. We apply \olup{height}-preserving
weakening (\olref[seq][adm]{prop:weak-G3c-adm}) to $\pi_1$ to obtain
\olup{proof} $\pi_1'$ of $B, \Gamma \fCenter \Delta', B \lif C, C,
A$ (we weaken with $B$ on the left and $C$ on the right).
Similarly, we apply \olref[seq][adm]{prop:weak-G3c-adm}
to~$\pi_3$ to get \olup{proof}~$\pi_3'$ of $A, B, \Gamma \fCenter
\Delta', B \lif C, C$ (by weakening with $B \lif C$ on the right).
The induction hypothesis applies to
\[
\AxiomC{}
\RightLabel{$\pi_1'$}
\Deduce$B, \Gamma \fCenter \Delta', B \lif C, C, A$
\AxiomC{}
\RightLabel{$\pi_3'$}
\Deduce$A, B, \Gamma \fCenter \Delta', B \lif C, C$
\RightLabel{\CutCS}
\BinaryInf$B, \Gamma \fCenter \Delta', B \lif C, C$
\DisplayProof
\]
with cut \olup{formula}~$A$: This \olup{proof} has the same cut rank and a
lower cut \olup{height} than~$\pi$. This gives us \olup{proof}~$\pi_4$ in
\Log{G3c} without \CutCS{} of~$B, \Gamma \fCenter \Delta', B \lif
C, C$. Apply $\RightR{\lif}$ to get \olup{proof} of $\Gamma \fCenter
\Delta', B \lif C, B \lif C$:
\[
\AxiomC{}
\RightLabel{$\pi_4$}
\Deduce$B, \Gamma \fCenter \Delta', B \lif C, C$
\RightLabel{$\RightR{\lif}$}
\UnaryInf$\Gamma \fCenter \Delta', B \lif C, B \lif C$
\DisplayProof
\]
To get \olup{proof}~$\pi'$ of $\Gamma \Sequent \Delta$, i.e., $\Gamma
\fCenter \Delta', B \lif C$, apply
\olref[seq][inv]{prop:G3c-cont-adm} (the admissibility of
contraction).

Now suppose $\pi$ ends in:
\[
\AxiomC{}
\RightLabel{$\pi_1$}
\Deduce$\lexists[x][B(x)], \Gamma' \fCenter \Delta, A$
\AxiomC{}
\RightLabel{$\pi_3$}
\Deduce$A, B(c), \Gamma' \fCenter \Delta$
\RightLabel{$\LeftR{\lexists}$}
\UnaryInf$A, \lexists[x][B(x)], \Gamma' \fCenter \Delta$
\RightSubproofLabel{$\pi_2$}
\RightLabel{\CutCS}
\BinaryInf$\Gamma' \fCenter \Delta$
\DisplayProof
\]
We again apply the inductive hypothesis to
\[
\AxiomC{}
\RightLabel{$\pi_1[B(c) \Sequent {}]$}
\Deduce$\lexists[x][B(x)], B(c), \Gamma' \fCenter \Delta, A$
\AxiomC{}
\LeftLabel{$\pi_3[\lexists[x][B(x)] \Sequent {}]$}
\Deduce$A, \lexists[x][B(x)], B(c), \Gamma' \fCenter \Delta$
\RightLabel{\CutCS}
\BinaryInf$\lexists[x][B(x)], B(c), \Gamma' \fCenter \Delta$
\RightLabel{\LeftR{\lexists}}
\UnaryInf$\lexists[x][B(x)], \lexists[x][B(x)], \Gamma' \fCenter \Delta$
\DisplayProof
\]
Note that the eigenvariable condition is satisfied, since the
eigenvariable condition on~\LeftR{\lexists} in~$\pi_2$ guarantees that
$c$ does not occur in $\lexists[x][B(x)]$, $\Gamma'$, or~$\Delta$.
Finally, apply \olref[seq][inv]{prop:G3c-cont-adm}.


\paragraph{Case D:}
At least one of $\pi_1$ and $\pi_2$ ends in a two-premise rule~$R$ in
which $A$ is not principal. There are 6 cases. We consider the case
where $A$ is not principal in the last inference of~$\pi_2$, which
ends in~$\RightR{\land}$; the other cases are analogous.

Assume $\pi$ ends in:
\[
\AxiomC{}
\RightLabel{$\pi_1$}
\Deduce$\Gamma \fCenter \Delta', B \land C, A$
\AxiomC{}
\RightLabel{$\pi_3$}
\Deduce$A, \Gamma \fCenter \Delta', B$
\AxiomC{}
\RightLabel{$\pi_4$}
\Deduce$A, \Gamma \fCenter \Delta', C$
\RightLabel{$\RightR{\land}$}
\BinaryInf$A, \Gamma \fCenter \Delta', B \land C$
\RightSubproofLabel{$\pi_2$}
\RightLabel{\CutCS}
\BinaryInf$\Gamma \fCenter \Delta$
\DisplayProof
\]
The inductive hypothesis applies to the \olup{proof}s
\[
\AxiomC{}
\RightLabel{$\pi_1'$}
\Deduce$\Gamma \fCenter \Delta', B \land C, B, A$
\AxiomC{}
\RightLabel{$\pi_3'$}
\Deduce$A, \Gamma \fCenter \Delta', B \land C, B$
\RightLabel{\CutCS}
\BinaryInf$\Gamma \fCenter \Delta', B \land C, B$
\DisplayProof
\]
and
\[
\AxiomC{}
\RightLabel{$\pi_1''$}
\Deduce$\Gamma \fCenter \Delta, B \land C, C, A$
\AxiomC{}
\RightLabel{$\pi_4'$}
\Deduce$A, \Gamma \fCenter \Delta', B \land C, C$
\RightLabel{\CutCS}
\BinaryInf$\Gamma \fCenter \Delta', B \land C, C$
\DisplayProof
\]
Here, the \Log{G3c} \olup{proof}s $\pi_1'$ and $\pi_1''$ are obtained by
\olup{height}-preserving weakening (\olref[seq][adm]{prop:weak-G3c-adm})
from~$\pi_1$ (once weakening on the right with $B$, and once with
$C$). The \olup{proof}s $\pi_3'$ and~$\pi_4'$ are obtained from $\pi_3$
and~$\pi_4$ also by \olup{height}-preserving weakening from $\pi_3$
and~$\pi_4$, respectively (by weakening with $B \land C$ on the
right).

The inductive hypothesis gives us \olup{proof}s~$\pi_5$ and~$\pi_6$ in
\Log{G3c} of $\Gamma \fCenter \Delta', B \land C, B$ and $\Gamma
\fCenter \Delta', B \land C, C$. We combine these into \olup{proof}
of $\Gamma \Sequent \Delta', B \land C, B \land C$ using the
rule~$\RightR{\land}$:
\[
\AxiomC{}
\RightLabel{$\pi_5$}
\Deduce$\Gamma \fCenter \Delta', B \land C, B$
\AxiomC{}
\RightLabel{$\pi_6$}
\Deduce$\Gamma \fCenter \Delta', B \land C, C$
\RightLabel{\RightR{\land}}
\BinaryInf$\Gamma \fCenter \Delta, B \land C, B \land C$
\DisplayProof
\]
To get \olup{proof} of $\Gamma \Sequent \Delta$, i.e., $\Gamma \fCenter
\Delta', B \land C$, apply \olref[seq][inv]{prop:G3c-cont-adm} (the
admissibility of contraction).

\paragraph{Case E: Reducing cuts.}
The final case concerns the situation where both $\pi_1$ and $\pi_2$
end in inferences in which $A$ is principal. There are six cases,
one for each possible \olup{main operator} of~$A$.

In each case, we turn a \CutCS{} with cut \olup{formula}~$A$ into one or more
cuts on sub-\olup{formula}s of~$A$, namely on the active \olup{formula}s of
the inferences of which~$A$ is the principal \olup{formula}. These cases
are therefore often referred to as ``reduction'' or ``conversion'' of
the cut inference.

\begin{enumerate}
\item If $A \ident \lnot B$ then $\pi$ ends in
\[
\AxiomC{}
\RightLabel{$\pi_1'$}
\Deduce$B, \Gamma \fCenter \Delta$
\RightLabel{\RightR{\lnot}}
\UnaryInf$\Gamma \fCenter \Delta, \lnot B$
\LeftSubproofLabel{$\pi_1$}
\AxiomC{}
\RightLabel{$\pi_2'$}
\Deduce$\Gamma \fCenter \Delta, B$
\RightLabel{\LeftR{\lnot}}
\UnaryInf$\lnot B, \Gamma \fCenter \Delta$
\RightSubproofLabel{$\pi_2$}
\RightLabel{\CutCS}
\BinaryInf$\Gamma \fCenter \Delta$
\DisplayProof
\]
By inductive hypothesis, \CutCS{} can be eliminated from
\[
\AxiomC{}
\RightLabel{$\pi_2'$}
\Deduce$\Gamma \fCenter \Delta, B$
\AxiomC{}
\RightLabel{$\pi_1'$}
\Deduce$B, \Gamma \fCenter \Delta$
\RightLabel{\CutCS}
\BinaryInf$\Gamma \fCenter \Delta$
\DisplayProof
\]

\item
If $A \ident (B \lor C)$ then $\pi$ ends in
\[
\AxiomC{}
\RightLabel{$\pi_1'$}
\Deduce$\Gamma \fCenter \Delta, B, C$
\RightLabel{\RightR{\lor}}
\UnaryInf$\Gamma \fCenter \Delta, B \lor C$
\LeftSubproofLabel{$\pi_1$}
\AxiomC{}
\RightLabel{$\pi_2'$}
\Deduce$B, \Gamma \fCenter \Delta$
\AxiomC{}
\RightLabel{$\pi_2''$}
\Deduce$C, \Gamma \fCenter \Delta$
\RightLabel{\LeftR{\lor}}
\BinaryInf$B \lor C, \Gamma \fCenter \Delta$
\RightSubproofLabel{$\pi_2$}
\RightLabel{\CutCS}
\BinaryInf$\Gamma \fCenter \Delta$
\DisplayProof
\]
Consider the \olup{proof} ending in \CutCS,
\[
\AxiomC{}
\RightLabel{$\pi_1'$}
\Deduce$\Gamma \fCenter \Delta, B, C$
\AxiomC{}
\RightLabel{$\pi_2'''$}
\Deduce$C, \Gamma \fCenter \Delta, B$
\RightLabel{\CutCS}
\BinaryInf$\Gamma \fCenter \Delta, B$
\DisplayProof
\]
where $\pi_2'''$ is obtained from $\pi_2''$ by \olup{height}-preserving weakening.
It has cut rank $< \cutr{\pi}$ as $\depth{C} < \depth{B \lor C}$, so the
inductive hypothesis yields \olup{proof} in \Log{G3c} of $\pi_1''$ of
$\Gamma \fCenter \Delta, B$. The \olup{proof} ending in \CutCS,
\[
\AxiomC{}
\RightLabel{$\pi_1''$}
\Deduce$\Gamma \fCenter \Delta, B$
\AxiomC{}
\RightLabel{$\pi_2'$}
\Deduce$B, \Gamma \fCenter \Delta$
\RightLabel{\CutCS}
\BinaryInf$\Gamma \fCenter \Delta$
\DisplayProof
\]
has cut rank $= \depth{B} < \depth{B \lor C}$, so the inductive
hypothesis applies again, and yields \olup{proof} $\pi'$ of $\Gamma
\fCenter \Delta$.

\item $A \ident (B \land C)$. Exercise.

\item If $A \ident (B \lif C)$ then $\pi$ ends in
\[
\AxiomC{}
\RightLabel{$\pi_1'$}
\Deduce$B, \Gamma \fCenter \Delta, C$
\RightLabel{\RightR{\lif}}
\UnaryInf$\Gamma \fCenter \Delta, B \lif C$
\LeftSubproofLabel{$\pi_1$}
\AxiomC{}
\RightLabel{$\pi_2'$}
\Deduce$\Gamma \fCenter \Delta, B$
\AxiomC{}
\RightLabel{$\pi_2''$}
\Deduce$C, \Gamma \fCenter \Delta$
\RightLabel{\LeftR{\lif}}
\BinaryInf$B \lif C, \Gamma \fCenter \Delta$
\RightSubproofLabel{$\pi_2$}
\RightLabel{\CutCS}
\BinaryInf$\Gamma \fCenter \Delta$
\DisplayProof
\]
The \olup{proof} ending in \CutCS,
\[
\AxiomC{}
\RightLabel{$\pi_1'$}
\Deduce$B, \Gamma \fCenter \Delta, C$
\AxiomC{}
\RightLabel{$\pi_2''$}
\Deduce$C, \Gamma \fCenter \Delta$
\RightLabel{\CutCS}
\BinaryInf$B, \Gamma \fCenter \Delta$
\DisplayProof
\]
has lower cut rank than~$\pi$. The inductive hypthesis yields
\olup{proof}~$\pi_1''$ of $B, \Gamma \fCenter \Delta$.
Now consider the \olup{proof} ending in \CutCS,
\[
\AxiomC{}
\RightLabel{$\pi_2'$}
\Deduce$\Gamma \fCenter \Delta, B$
\AxiomC{}
\RightLabel{$\pi_1''$}
\Deduce$B, \Gamma \fCenter \Delta$
\RightLabel{\CutCS}
\BinaryInf$\Gamma \fCenter \Delta$
\DisplayProof
\]
Its cut rank is also less than the cut rank of~$\pi$, so the induction
hypothesis yields a \Log{G3c}-\olup{proof}~$\pi'$ of $\Gamma \fCenter
\Delta$.
\item If $A \ident \lforall[x][B(x)]$ then $\pi$ ends in
\[
\AxiomC{}
\RightLabel{$\pi_1'$}
\Deduce$\Gamma \fCenter \Delta, B(c)$
\RightLabel{\RightR{\lforall}}
\UnaryInf$\Gamma \fCenter \Delta, \lforall[x][B(x)]$
\LeftSubproofLabel{$\pi_1$}
\AxiomC{}
\RightLabel{$\pi_2'$}
\Deduce$B(t), \lforall[x][B(x)], \Gamma \fCenter \Delta$
\RightLabel{\LeftR{\lforall}}
\UnaryInf$\lforall[x][B(x)], \Gamma \fCenter \Delta$
\RightSubproofLabel{$\pi_2$}
\RightLabel{\CutCS}
\BinaryInf$\Gamma \fCenter \Delta$
\DisplayProof
\]
By inductive hypothesis, \CutCS{} can be eliminated from
\[
\AxiomC{}
\RightLabel{$\pi_1''$}
\Deduce$B(t), \Gamma \fCenter \Delta, \lforall[x][B(x)]$
\AxiomC{}
\RightLabel{$\pi_2'$}
\Deduce$B(t), \lforall[x][B(x)], \Gamma \fCenter \Delta$
\RightLabel{\CutCS}
\BinaryInf$B(t), \Gamma \fCenter \Delta$
\DisplayProof
\]
where $\pi_1''$ is obtained from $\pi_1$ (not~$\pi_1'$!) by
\olup{height}-preserving weakening. (The \olup{proof} has the same cut rank
but lower cut \olup{height}.) This gives us \olup{proof}~$\pi_2''$ of
$B(t), \Gamma \fCenter \Delta$.

The \olup{proof}~$\pi_1'$ has end-sequent $\Gamma \Sequent \Delta,
B(c)$, where $c$~is the eigenvariable of a~$\RightR{\lforall}$
inference. So $c$~does not occur in~$B(x)$, $\Gamma$, or~$\Delta$. We
assume the \olup{proof}~$\pi$ is regular, so $c$~is the eigenvariable of
only the following~$\RightR{\lforall}$ inference and only occurs above
it, i.e., only occurs in~$\pi_1'$ and is not an eigenvariable of an
inference in~$\pi_1'$. Since $c$ does not occur in~$\pi_2$, it cannot
occur in~$t$. By \olref[seq][qua]{lem:subst-regular}, the
\olup{proof}~$\Subst{\pi_1'}{t}{c}$ is \olup{proof} of $\Gamma \Sequent \Delta,
B(t)$. We now apply the inductive hypothesis to the \olup{proof}
\[
\AxiomC{}
\RightLabel{$\Subst{\pi_1'}{t}{c}$}
\Deduce$\Gamma \fCenter \Delta, B(t)$
\AxiomC{}
\RightLabel{$\pi_2''$}
\Deduce$B(t), \Gamma \fCenter \Delta$
\RightLabel{\CutCS}
\BinaryInf$\Gamma \fCenter \Delta$
\DisplayProof
\]
(The induction hypothesis applies since the cut \olup{formula}
is~$B(t)$, so the \olup{proof} has lower cut rank than~$\pi$.)
\item We leave the case $A \ident \lexists[x][B(x)]$ as an exercise.
\end{enumerate}
\end{proof}

\begin{prob}
Verify the sub-case of case~D in the proof of \olref{lem:cut-adm-G3c}
where $A$ is not principal in the last inference of~$\pi_1$, which
ends in~$\LeftR{\lif}$. (Note: part of this exercise is to clearly
state what you have to prove, i.e., what form $\pi$ has in this case.)
\end{prob}

\begin{prob}
Verify the sub-case of case~D in the proof of \olref{lem:cut-adm-G3c}
where $A$ is not principal in the last inference of~$\pi_1$, which
ends in~$\LeftR{\lforall}$.
\end{prob}

\begin{prob}
Verify the sub-case of case~E in the proof of \olref{lem:cut-adm-G3c}
where $A$ is principal in the last inferences of both $\pi_1$ and
$\pi_2$ and $A \ident (B \land C)$.
\end{prob}

\begin{prob}
Verify the sub-case of case~E in the proof of \olref{lem:cut-adm-G3c}
where $A$ is principal in the last inferences of both $\pi_1$ and
$\pi_2$ and $A \ident \lexists[x][B(x)]$.
\end{prob}

Note that in case~E the cut \olup{height} of the \olup{proof} ending in
\CutCS{} to which we apply the inductive hypothesis can be larger than
the cut \olup{height} of the original proof. E.g., in the case where $A
\ident (B \lif C)$ we turned $\pi$ into \olup{proof} involving two
\CutCS{} inferences,
\[
\AxiomC{}
\RightLabel{$\pi_2'$}
\Deduce$\Gamma \fCenter \Delta, B$
\AxiomC{}
\RightLabel{$\pi_1'$}
\Deduce$B, \Gamma \fCenter \Delta, C$
\AxiomC{}
\RightLabel{$\pi_2''$}
\Deduce$C, \Gamma \fCenter \Delta$
\RightLabel{\CutCS}
\BinaryInf$B, \Gamma \fCenter \Delta$
\RightSubproofLabel{$\pi_1''$}
\RightLabel{\CutCS}
\BinaryInf$B, \Gamma \fCenter \Delta$
\DisplayProof
\]
The top \CutCS{} between $\pi_1'$ and $\pi_2''$ has both lower cut
rank and lower cut \olup{height} than~$\pi$. But the
\olup{proof}~$\pi_1''$---the result of applying the inductive
hypothesis---no longer has lower cut \olup{height} than~$\pi$. Since the
cut \olup{formula} of the lower \CutCS{} is $B$, its cut \emph{rank},
however, is lower than the cut rank of~$\pi$.

\olref{lem:cut-adm-G3c} establishes that we can eliminate a single topmost
\CutCS{} inference from \olup{proof}. We can show that we can eliminate
any number of \CutCS{} inferences from \olup{proof} by successively
applying it to the topmost sub-\olup{proof}s ending in~\CutCS.

\begin{thm}\ollabel{thm:cut-elim-G3c-topmost} Any \olup{proof} $\pi$ in
$\Log{G3c} + \CutCS$ can be transformed into \olup{proof} in~\Log{G3c}.
\end{thm}

\begin{proof}
  By induction on the number~$n$ of \CutCS{} inferences in~$\pi$. If
  $n=0$, there is nothing to do: $\pi$ is already \olup{proof}
  in~\Log{G3c}. Otherwise, consider the sub-\olup{proof}~$\pi'$ ending in
  a top-most \CutCS{} inference. It contains a single \CutCS{} rule as the
  last inference. Apply \olref{lem:cut-adm-G3c} to transform it into a
  \CutCS{}-free proof~$\pi''$, and replace the sub-\olup{proof}~$\pi'$
  by~$\pi''$. The result is \olup{proof} containing $n-1$ \CutCS{}
  inferences, and the induction hypothesis applies.
\end{proof}

% ==== OLP file: content/proof-theory/cut-elimination/cut-elimination.tex ====
% Part: proof-theory
% Chapter: cut-elimination



\olchapter{pt}{cut}{Cut Elimination}

\olimport{introduction}
\olimport{ce-topmost}
\olimport{ce-largest}
\olimport{midsequent}
\olimport{interpolation}

\OLEndChapterHook


% ==== OLP file: content/proof-theory/cut-elimination/interpolation.tex ====
% Part: proof-theory
% Chapter: cut-elimination
% Section: interpolation



\olfileid{pt}{cut}{itp}

\olsection{Maehara's Lemma and Craig's Interpolation Theorem}

The interpolation theorem, due to William Craig, states that if $A
\Entails B$ then there is \olup{sentence}~$C$ (the \emph{interpolant})
such that $A \Entails C$ and $C \Entails B$. Importantly, the
interpolant~$C$ can be chosen so that the \olup{constant}s and
\olup{predicate}s occurring in~$C$ all occur in both~$A$ and~$B$. (We
assume that no \olup{function}s are involved.) Without this restriction,
$A$~and $B$ could trivially serve as interpolants.

The interpolation theorem holds for classical first-order logic and
has been called the ``last important property of first-order logic''
to be discovered. It has important applications in model theory, and
automated reasoning. Its original motivation and application was in
the philosophy of science, since it can be used to study the
primitives by which a theory can or cannot be axiomatized. It also
implies Beth's definability theorem, which states that if a theory can
implicitly define a concept, it can define it explicitly.

Like many results in mathematical logic, the interpolation theorem can
be proved both using the methods of model theory and those of proof
theory. The advantage of proof theoretic methods is that not only do
they show the existence of an interpolant, but provide an algorithm
that computes the interpolant from \olup{proof} that $A \Entails B$,
e.g., \olup{proof} in~\Log{G3c} of $A \Sequent B$.

Let $\Lang L(A)$ be the set of \olup{constant}s and \olup{predicate}s
in~$A$ and $\Lang L(\Gamma) = \bigcup \Setabs{\Lang L(A)}{A \in
\Gamma}$. Throughout this section, we assume we are dealing with
\olup{formula}s without \olup{function}s.

\begin{thm}[Craig's Interpolation Theorem]\ollabel{thm:interpolation}
Let the \olup{language} of~$A$ be $\Lang L(A)$ and that of $B$ be
$\Lang L(B)$. If $\Log{G3c} \Proves A \Sequent B$ then there is
\olup{formula}~$C$ in the \olup{language}~$\Lang L(C) \subseteq \Lang
L(A) \cap \Lang L(B)$ such that $\Log{G3c} \Proves A \Sequent C$
and $\Log{G3c} \Proves C \Sequent B$.
\end{thm}

We prove a generalization of the interpolation theorem to make use of
the structure of \olup{proof}s in~\Log{G3c}. To this end, imagine the
antecendent and succedents of sequents to be separated into two parts
each. We write $\maeh{\Gamma_1}{\Gamma_2} \Sequent
\maeh{\Delta_1}{\Delta_2}$ to indicate that $\Gamma_1$ and $\Delta_1$
belong to the first part and $\Gamma_2$ and $\Delta_2$ to the second
part. The interpolation theorem then follows immediately from the
following lemma.

\begin{lem}[Maehara's Lemma]\ollabel{lem:maehara} If $\Log{G3c}
\Proves \maeh{\Gamma_1}{\Gamma_2} \Sequent \maeh{\Delta_1}{\Delta_2}$
then there is \olup{formula}~$C$ in the \olup{language}~$\Lang L(C)
\subseteq \Lang L(\Gamma_1,
\Delta_1) \cap \Lang L(\Gamma_2, \Delta_2)$ such that $\Log{G3c}
\Proves \Gamma_1 \Sequent \Delta_1, C$ and $\Log{G3c} \Proves C,
\Gamma_2 \Sequent \Delta_2$.
\end{lem}

\begin{proof}
  Assume that $\pi \Proves \maeh{\Gamma_1}{\Gamma_2} \Sequent
  \maeh{\Delta_1}{\Delta_2}$. We prove the claim by induction on
  $n=\pheight{\pi}$.

  If $n=0$, then $\maeh{\Gamma_1}{\Gamma_2} \Sequent
  \maeh{\Delta_1}{\Delta_2}$ is an axiom, and some atomic
  \olup{formula}~$A$ occurs in both $\Gamma_1, \Gamma_2$ and $\Delta_1,
  \Delta_2$. There are four cases depending on whether $A$~belongs to
  the first part or the second on the left, and whether it belongs to
  the first or second part on the right.
  \begin{enumerate}
    \item $A \in \Gamma_1$ and $A \in \Delta_1$. We want a~$C$ such
    that both $\Gamma_1 \Sequent \Delta_1, C$ and $C, \Gamma_2
    \Sequent \Delta_2$ are \olup{provable}. The first is already an axiom
    regardless of how we pick~$C$ since $A$ occurs on both the left
    and right. Our only option to guarantee that $C, \Gamma_2
    \Sequent \Delta_2$ is \olup{provable} is to let $C \ident \lfalse$.
    \item $A \in \Gamma_1$ and $A \in \Delta_2$. Then $\Gamma_1
    \Sequent \Delta_1, A$ and $A, \Gamma_2 \Sequent \Delta_2$ are
    both axioms, so we can pick $C \ident A$.
    \item $A \in \Gamma_2$ and $A \in \Delta_1$. Then $A, \Gamma_1
    \Sequent \Delta_1$ and $\Gamma_2 \Sequent \Delta_2, A$ would be 
    axioms, but they have~$A$ on the wrong sides. So $A$ isn't the
    interpolant. However, using $\RightR{\lnot}$ and $\LeftR{\lnot}$
    we can \olup{prove} $\Gamma_1
    \Sequent \Delta_1, \lnot A$ and $\lnot A, \Gamma_2 \Sequent
    \Delta_2$.
    \item $A \in \Gamma_2$ and $A \in \Delta_2$. Exercise.
  \end{enumerate}
  The cases where $\maeh{\Gamma_1}{\Gamma_2} \Sequent
  \maeh{\Delta_1}{\Delta_2}$ is an axiom because $\lfalse \in
  \Gamma_1$ or $\lfalse \in \Gamma_2$ are left as exercises.

  If $n>0$, $\pi$ ends in a logical inference. We have to distinguish
  cases cases according to which rule was used, and to which part the
  principal formula belongs. In each case we can use the induction
  hypothesis that the lemma holds for all \olup{proof}s~$\pi_1$ with
  $\pheight{\pi_1} < n$, and however the end-sequent of~$\pi_1$ may be
  partitioned into two parts. Let's consider some specific cases.
  \begin{enumerate}
    \item $\pi$ ends in $\RightR{\lnot}$ with principal formula~$\lnot
    A$. We have two cases: depending on whether $\lnot A$ belongs to
    the first or second part, $\pi$ ends in either
    \[
    \AxiomC{}
    \RightLabel{$\pi_1$}
    \Deduce$\maeh{A, \Gamma_1}{\Gamma_2} \fCenter
      \maeh{\Delta_1}{\Delta_2}$
    \RightLabel{\RightR{\lnot}}
    \UnaryInf$\maeh{\Gamma_1}{\Gamma_2} \fCenter
      \maeh{\Delta_1, \lnot A}{\Delta_2}$
    \DisplayProof
    \]
    or
    \[
    \AxiomC{}
    \RightLabel{$\pi_1$}
    \Deduce$\maeh{\Gamma_1}{A, \Gamma_2} \fCenter
      \maeh{\Delta_1}{\Delta_2}$
    \RightLabel{\RightR{\lnot}}
    \UnaryInf$\maeh{\Gamma_1}{\Gamma_2} \fCenter
      \maeh{\Delta_1}{\Delta_2, \lnot A}$
    \DisplayProof
    \]
    Note that if the principal \olup{formula}~$\lnot A$ belongs to the
    first part we assigned the active \olup{formula}~$A$ also to the
    first part the premise. This is a choice we get to make. The
    inductive hypothesis would also apply if we assigned the active
    formula to the other part, but then we would not be able to find
    an interpolant (exercise: try it).

    Consider the first situation first, where $\lnot A$~belongs
    to the first part. By inductive hypothesis, there is an
    interpolant of the end-sequent of~$\pi_1$, i.e., \olup{formula}~$C_1$
    such that both
    \begin{align*}
    A, \Gamma_1 & \Sequent \Delta_1, C_1 &&\text{and} &
    C_1, \Gamma_2 & \Sequent \Delta_2
    \intertext{are provable. What we need is \olup{formula}~$C$ such that
    both of the following are \olup{provable}:}
    \Gamma_1 & \Sequent \Delta_1, \lnot A, C &&\text{and} &
    C, \Gamma_2 & \Sequent \Delta_2
    \end{align*}
    Clearly, we can just take $C \ident C_1$: the inductive hypothesis
    already told us that the sequent on the right is \olup{provable}, and
    the sequent on the left follows from the sequent we got from the
    inductive hypothesis by~$\RightR{\lnot}$. The inductive hypothesis
    also tells us that $\Lang L(C_1) \subseteq \Lang L(A, \Gamma_1,
    \Delta_1) \cap \Lang L(\Gamma_2, \Delta_2)$. Since $\Lang L(\lnot
    A) = \Lang L(A)$ and $\Lang L(C) = \Lang L(C_1)$ we have that
    $\Lang L(C_1) \subseteq \Lang L(\Gamma_1, \Delta_1, \lnot A) \cap \Lang
    L(\Gamma_2, \Delta_2)$.

    In the second scenario the situation is slightly different
    in that $\lnot A$~belongs to the second part. We assign the
    active \olup{formula} in the premise also to the second part and
    apply the inductive hypothesis to
    \begin{align*}
    \Gamma_1 & \Sequent \Delta_1, C_1 &&\text{and} &
    C_1, A, \Gamma_2 & \Sequent \Delta_2
    \intertext{are provable. If we again apply the \RightR{\lnot} rule to the second of these, we get \olup{proof}s of}
    \Gamma_1 & \Sequent \Delta_1, C_1 &&\text{and} &
    C_1, \Gamma_2 & \Sequent \Delta_2, \lnot A
    \end{align*}
    These are again exactly the sequents we want to be \olup{provable} if
    $C_1$ is an interpolant; we again let $C \ident C_1$. $\Lang
    L(C_1) \subseteq \Lang L(\Gamma_1, \Delta_1) \cap \Lang L(\Gamma_2,
    \Delta_2, \lnot A)$ as before.
    \item $\pi$ ends in $\RightR{\lif}$ with principal \olup{formula}~$A
    \lif B$. We again have two cases, depending on whether $A \lif
    B$ belongs to the first or second part. The \olup{proof}~$\pi$ ends
    in either
    \[
    \AxiomC{}
    \RightLabel{$\pi_1$}
    \Deduce$\maeh{A, \Gamma_1}{\Gamma_2} \fCenter
      \maeh{\Delta_1, B}{\Delta_2}$
    \RightLabel{\RightR{\lif}}
    \UnaryInf$\maeh{\Gamma_1}{\Gamma_2} \fCenter
      \maeh{\Delta_1, A \lif B}{\Delta_2}$
    \DisplayProof
    \]
    or
    \[
    \AxiomC{}
    \RightLabel{$\pi_1$}
    \Deduce$\maeh{\Gamma_1}{A, \Gamma_2} \fCenter
      \maeh{\Delta_1}{\Delta_2, B}$
    \RightLabel{\RightR{\lif}}
    \UnaryInf$\maeh{\Gamma_1}{\Gamma_2} \fCenter
      \maeh{\Delta_1}{\Delta_2, A \lif B}$
    \DisplayProof
    \]
    Again we have assigned the active \olup{formula}s to the same part in
    the premise as the principal \olup{formula} in the conclusion. In the
    first secenario, we apply the inductive hypothesis and get
    \olup{proof}s of
    \begin{align*}
    A, \Gamma_1 & \Sequent \Delta_1, B, C_1 &&\text{and} & C_1, \Gamma_2 & \Sequent \Delta_2
    \intertext{The sequent on the left is a suitable premise for \RightR{\lif}. So we have the following \olup{provable} sequents:}
    \Gamma_1 & \Sequent \Delta_1, A \lif B, C_1 &&\text{and} & C_1, \Gamma_2 & \Sequent \Delta_2
    \end{align*}
    This means that $C_1$ is also an interpolant for the end-sequent
    of~$\pi$ and we can take $C \ident C_1$ once more.
    
    The other scenario is handled the same way.
    \item $\pi$ ends in $\LeftR{\lif}$ with principal \olup{formula}~$A
    \lif B$. If $A \lif B$ belongs to the first part, the \olup{proof}~$\pi$ ends
    in
    \[
    \AxiomC{}
    \RightLabel{$\pi_1$}
    \Deduce$\maeh{\Gamma_1}{\Gamma_2} \fCenter
      \maeh{\Delta_1, A}{\Delta_2}$
    \AxiomC{}
    \RightLabel{$\pi_2$}
    \Deduce$\maeh{B, \Gamma_1}{\Gamma_2} \fCenter
      \maeh{\Delta_1}{\Delta_2}$
    \BinaryInf$\maeh{A \lif B, \Gamma_1}{\Gamma_2} \fCenter
      \maeh{\Delta_1}{\Delta_2}$
    \DisplayProof
    \]
    The inductive hypothesis now gives us four \olup{provable} sequents,
    two for the interpolant $C_1$ of the left premise and two for the
    interpolant $C_2$ of the second premise:
    \begin{align*}
    \Gamma_1 & \Sequent \Delta_1, A, C_1 &&\text{(1)} &
    C_1, \Gamma_2 & \Sequent \Delta_2 && \text{(2)}\\
    B, \Gamma_1 & \Sequent \Delta_1, C_2 &&\text{(3)} &
    C_2, \Gamma_2 & \Sequent \Delta_2 && \text{(4)}
    \end{align*}
    Sequents (1) and (3) can serve as premises of~\LeftR{\lif}, if we
    apply weakening and add $C_2$ on the right to~(1) and $C_1$ on
    the right to~(3):
    \[
    \AxiomC{}
    \Deduce$\Gamma_1 \fCenter \Delta_1, A, C_1, C_2$
    \AxiomC{}
    \Deduce$B, \Gamma_1 \fCenter \Delta_1, C_1, C_2$
    \RightLabel{\LeftR{\lif}}
    \BinaryInf$A \lif B, \Gamma_1 \fCenter \Delta_1, C_1, C_2$
    \DisplayProof
    \]
    By applying $\RightR{\lor}$ we get the sequent
    \[
    A \lif B, \Gamma_1 \fCenter \Delta_1, C_1 \lor C_2
    \]
    This suggests that $C_1 \lor C_2$ might be our interpolant in
    this case. And indeed, we can \olup{prove} the other required sequent
    from the remaining sequents (2) and~(4):
    \[
    \AxiomC{}
    \Deduce$C_1, \Gamma_2 \fCenter \Delta_2$
    \AxiomC{}
    \Deduce$C_2, \Gamma_2 \fCenter \Delta_2$
    \RightLabel{\LeftR{\lor}}
    \BinaryInf$C_1 \lor C_2, \Gamma_2 \fCenter \Delta_2$
    \DisplayProof
    \]
    It remains to verify that $C_1 \lor C_2$ is in the correct
    language. The inductive hypothesis gives us that 
    \begin{align*}
    \Lang L(C_1) & \subseteq
    \Lang L(\Gamma_1, \Delta_1, A) \cap \Lang L(\Gamma_2, \Delta_2) &\text{and}\\
    \Lang L(C_2) & \subseteq
    \Lang L(B, \Gamma_1, \Delta_1) \cap \Lang L(\Gamma_2, \Delta_2)
    \intertext{Since}
    \Lang L(A) & \subseteq \Lang L(A \lif B) &\text{and}\\
    \Lang L(B) & \subseteq \Lang L(A \lif B)
    \intertext{we have both}
    \Lang L(C_1) & \subseteq
    \Lang L(\Gamma_1, \Delta_1, A \lif B) \cap \Lang L(\Gamma_2, \Delta_2) &\text{and}\\
    \Lang L(C_2) & \subseteq
    \Lang L(\Gamma_1, \Delta_1, A \lif B) \cap \Lang L(\Gamma_2, \Delta_2)
    \intertext{and consequently, $\Lang L (C_1 \lor C_2) = {}$}
    \Lang L(C_1) \cup \Lang L(C_2) & \subseteq
    \Lang L(\Gamma_1, \Delta_1, A \lif B) \cap \Lang L(\Gamma_2, \Delta_2),
    \end{align*}
    which is as required.

    In the case where $A \lif B$ belongs to the second part, the
    situation is different, but handled in a similar way. The
    inductive hypothesis again gives us four \olup{provable} sequents:
    \begin{align*}
    \Gamma_1 & \Sequent \Delta_1, C_1 &&\text{(1)} &
    C_1, \Gamma_2 & \Sequent \Delta_2, A && \text{(2)}\\
    \Gamma_1 & \Sequent \Delta_1, C_2 &&\text{(3)} &
    C_2, B, \Gamma_2 & \Sequent \Delta_2 && \text{(4)}
    \end{align*}
    Now sequents (2) and (4)---plus some weakened \olup{formula}s---serve as premises of~\LeftR{\lif}:
    \[
    \AxiomC{}
    \Deduce$C_1, C_2, \Gamma_2 \fCenter \Delta_2, A$
    \AxiomC{}
    \Deduce$B, C_1, C_2, \Gamma_2 \fCenter \Delta_2$
    \RightLabel{\LeftR{\lif}}
    \BinaryInf$A \lif B, C_1, C_2, \Gamma_2 \fCenter \Delta_2$
    \DisplayProof
    \]
    We again want to turn this into a sequent that has a single
    \olup{formula} built from $C_1$ and $C_2$, but this time we need it
    in the antecedent (since we're now dealing with the second part).
    Applying $\LeftR{\land}$ does the trick; we should take $C \ident
    (C_1 \land C_2)$ as the interpolant. Fortuitously, we can use
    the remaining sequents (1) and (3) to \olup{prove} the corresponding
    sequent for the other part:
    \[
    \AxiomC{}
    \Deduce$\Gamma_1 \fCenter \Delta_1, C_1$
    \AxiomC{}
    \Deduce$\Gamma_1 \fCenter \Delta_1, C_2$
    \RightLabel{\RightR{\land}}
    \BinaryInf$\Gamma_1 \fCenter \Delta_1, C_1 \land C_2$
    \DisplayProof
    \]
    As before, we have $\Lang L(C_1 \land C_2) \subseteq \Lang
    L(\Gamma_1, \Delta_1) \cap \Lang L(\Gamma_2, \Delta_2, A \lif
    B)$.

    \item $\pi$ ends in $\RightR{\lforall}$ with principal \olup{formula}~$\lforall[x][A(x)]$:
    \[
    \AxiomC{}
    \RightLabel{$\pi_1$}
    \Deduce$\maeh{\Gamma_1}{\Gamma_2} \fCenter
    \maeh{\Delta_1, A(c)}{\Delta_2}$
    \RightLabel{\RightR{\lforall}}
    \UnaryInf$\maeh{\Gamma_1}{\Gamma_2} \fCenter
      \maeh{\Delta_1, \lforall[x][A(x)]}{\Delta_2}$
    \DisplayProof
    \]
    or
    \[
    \AxiomC{}
    \RightLabel{$\pi_1$}
    \Deduce$\maeh{\Gamma_1}{\Gamma_2} \fCenter
      \maeh{\Delta_1}{\Delta_2, A(c)}$
    \RightLabel{\RightR{\lforall}}
    \UnaryInf$\maeh{\Gamma_1}{\Gamma_2} \fCenter
      \maeh{\Delta_1}{\Delta_2, \lforall[x][A(x)]}$
    \DisplayProof
    \]
    In the first scenario, the induction hypothesis gives us \olup{proof}s
    of $\Gamma_1 \Sequent \Delta_1, A(c), C_1$ and $C_1, \Gamma_2
    \Sequent \Delta_2$. We can apply $\RightR{\lforall}$ to the first
    to obtain $\Gamma_1 \Sequent \Delta_1, \lforall[x][A(x)], C_1$.
    (The eigenvariable condition is satisfied, since it was already
    satisfied in the \RightR{\lforall} rule at the end of~$\pi$.) Thus
    $C_1$ would be an interpolant if the condition that $\Lang
    L(C_1) \subseteq \Lang L(\Gamma_1, \Delta_1, \lforall[x][A(x)])
    \cap \Lang L(\Gamma_2, \Delta_2)$ is satisfied. By the inductive
    hypothesis, we have $\Lang L(C_1) \subseteq \Lang L(\Gamma_1,
    \Delta_1, A(c)) \cap \Lang L(\Gamma_2, \Delta_2)$. So we have to
    worry about~$c$---can $c$ occur in~$C_1$? It can't. If $c$ occurs
    in $C_1$, we would have both $c \in \Lang L(\Gamma_1, \Delta_1,
    A(c))$ (which is true) and also $c \in \Lang L(\Gamma_2,
    \Delta_2)$. But then $c$ would occur in the conclusion of the
    $\RightR{\lforall}$ inference in~$\pi$, violating the eigenvariable
    condition.
    
    The other scenario is similar.

    \item $\pi$ ends in $\RightR{\lexists}$ with principal
    \olup{formula}~$\lexists[x][A(x)]$. We again have two scenarios. We
    consider the scenario where $\lexists[x][A(x)]$ belongs to the
    first part, and leave the other as an exercise.

    In the first scenario, $\pi$ ends in
    \[
    \AxiomC{}
    \RightLabel{$\pi_1$}
    \Deduce$\maeh{\Gamma_1}{\Gamma_2} \fCenter
      \maeh{\Delta_1, \lexists[x][A(x)], A(t)}{\Delta_2}$
    \RightLabel{\RightR{\lexists}}
    \UnaryInf$\maeh{\Gamma_1}{\Gamma_2} \fCenter
      \maeh{\Delta_1, \lexists[x][A(x)]}{\Delta_2}$
    \DisplayProof
    \]
    The induction hypothesis gives us \olup{proof}s of $\Gamma_1 \Sequent
    \Delta_1, \lexists[x][A(x)], A(t), C_1$ and $C_1, \Gamma_2
    \Sequent \Delta_2$. We can apply $\RightR{\lexists}$ to the first
    to obtain $\Gamma_1 \Sequent \Delta_1, \lexists[x][A(x)], C_1$.
    Thus $C_1$ would again be an interpolant if the condition that
    $\Lang L(C_1) \subseteq \Lang L(\Gamma_1, \Delta_1,
    \lexists[x][A(x)]) \cap \Lang L(\Gamma_2, \Delta_2)$ is
    satisfied. Here the inductive hypothesis gives us $\Lang L(C_1)
    \subseteq \Lang L(\Gamma_1, \Delta_1, \lexists[x][A(x)], A(t))
    \cap \Lang L(\Gamma_2, \Delta_2)$. Now remember that we assumed
    that there are no \olup{function}s. That means the term~$t$ can only
    be \olup{constant}, i.e., $t \ident b$. If $b \notin \Lang L(C_1)$
    or $b \in \Lang L(\Gamma_1, \Delta_1, \lexists[x][A(x)]) \cap
    \Lang L(\Gamma_2, \Delta_2)$ there is nothing to worry about: we
    can take $C_1$ as the interpolant. But what if $b \in \Lang
    L(C_1)$ and $b \notin \Lang L(\Gamma_1, \Delta_1,
    \lexists[x][A(x)]) \cap \Lang L(\Gamma_2, \Delta_2)$? First of
    all, by replacing all occurrences of~$b$ in~$C_1$ by
    \olup{variable}~$y$ we can write $C_1 \ident \Subst{C_1'}{b}{y}$
    where $C_1'$~does not contain~$b$. Furthermore, either $b \notin
    \Lang L(\Gamma_1, \Delta_1, \lexists[x][A(x)])$ or $b \notin
    \Lang L(\Gamma_2, \Delta_2)$. We can take $C \ident
    \lforall[y][C_1'(y)]$ or $C \ident \lexists[y][C_1'(y)]$,
    respectively. In the first case, we have:
    \[
    \AxiomC{}
    \Deduce$\Gamma_1 \fCenter
      \Delta_1, \lexists[x][A(x)], A(b), C_1'(b)$
    \RightLabel{\RightR{\lexists}}
    \UnaryInf$\Gamma_1 \fCenter \Delta_1, \lexists[x][A(x)], C_1'(b)$
    \RightLabel{\RightR{\lforall}}
    \UnaryInf$\Gamma_1 \fCenter
      \Delta_1, \lexists[x][A(x)], \lforall[y][C_1'(y)]$
    \DisplayProof
    \]
    and
    \[
    \AxiomC{}
    \Deduce$C_1'(b), \lforall[y][C_1'(y)], \Gamma_2 \fCenter \Delta_2$
    \RightLabel{\LeftR{\lforall}}
    \UnaryInf$\lforall[y][C_1'(y)], \Gamma_2 \fCenter \Delta_2$
    \DisplayProof
    \]
    The \olup{proof}s of the top sequents are from the inductive
    hypothesis (in the second one, we apply admissibility of weakening
    with $\lforall[y][C_1'(y)]$ in the antecedent). The eigenvariable
    condition on \RightR{\lforall} in the first \olup{proof} is satisfied
    because $b \notin \Lang L(\Gamma_1, \Delta_1,
    \lexists[x][A(x)])$, and $C_1'$~was defined so it does not
    contain~$b$.
    
    In the second case, we have (from the inductive hypothesis plus
    admissibility of weakening by~$\lexists[y][C_1'(y)]$):
    \[
    \AxiomC{}
    \Deduce$\Gamma_1 \fCenter
      \Delta_1, \lexists[x][A(x)], A(b), \lexists[y][C_1'(y)], C_1'(b)$
    \RightLabel{\RightR{\lexists}}
    \UnaryInf$\Gamma_1 \fCenter
      \Delta_1, \lexists[x][A(x)], \lexists[y][C_1'(y)], C_1'(b)$
    \RightLabel{\RightR{\lexists}}
    \UnaryInf$\Gamma_1 \fCenter
      \Delta_1, \lexists[x][A(x)], \lexists[y][C_1'(y)]$
    \DisplayProof
    \]
    and
    \[
    \AxiomC{}
    \Deduce$C_1'(b), \Gamma_2 \fCenter \Delta_2$
    \RightLabel{\LeftR{\lexists}}
    \UnaryInf$\lexists[y][C_1'(y)], \Gamma_2 \fCenter \Delta_2$
    \DisplayProof
    \]
    The eigenvariable condition on \LeftR{\lexists} in the second
    \olup{proof} is satisfied because $b \notin \Lang L(\Gamma_2,
    \Delta_2)$, and $C_1'$ was defined so it does not contain~$b$.
  \end{enumerate}
  The other cases are similary; we leave them as exercises.
\end{proof}

\begin{prob}
Suppose we had a connective $\lnand$ with rules
\[
\Axiom$\Gamma \fCenter \Delta, A$
\Axiom$\Gamma \fCenter \Delta, B$
\RightLabel{\LeftR{\lnand}}
\BinaryInf$A \lnand B, \Gamma \fCenter \Delta$
\DisplayProof
\quad
\Axiom$A, B, \Gamma \fCenter \Delta$
\RightLabel{\RightR{\lnand}}
\UnaryInf$\Gamma \fCenter \Delta, A \lnand B$
\DisplayProof
\]
Show that Maehara's Lemma still holds by carrying out the cases in the
proof of \olref{lem:maehara} where $\pi$ ends in one of these rules.
\end{prob}

Having established Maehara's Lemma, we can use it to prove Craig's
Interpolation Theorem.

\begin{proof}[Proof of \olref{thm:interpolation}]
  Suppose $A \Sequent B$ is \olup{provable}. Apply \olref{lem:maehara}
  to $\maeh{A}{\ } \Sequent \maeh{\ }{B}$ to obtain the
  interpolant~$C$. We have $\Proves A \Sequent C$ and $\Proves C
  \Sequent B$.
\end{proof}

Suppose $\Gamma$ is a theory (set of \olup{sentence}s) in
\olup{language}~$\Lang L$ containing an $n$-place {predicate}~$R$. We
say that $\Gamma$ \emph{implicitly defines}~$R$, iff
\begin{align*}
\Gamma, \Gamma'
& \Entails \lforall[x_1][\dots\lforall[x_n][(R(x_1,\dots,x_n) \liff
R'(x_1, \dots, x_n))]],
\intertext{where $\Gamma'$ is $\Gamma$ with all
occurrences of~$R$ replaced by~$R' \notin \Lang L$. $\Gamma$
\emph{explicitly defines}~$R$ iff there is some~$A(x_1, \dots, x_n)$
not containing~$R$ such that}
\Gamma &\Entails
\lforall[x_1][\dots\lforall[x_n][(R(x_1,\dots,x_n) \liff A(x_1,
\dots, x_n))]].
\end{align*}
Clearly, if $\Gamma$ explicitly defines $R$ it implicitly defines~$R$.
The converse is not obvious, but nevertheless holds:

\begin{lem}[Beth Definability Theorem]
  If $\Gamma$ implicitly defines~$R$ then it also explicitly defines it.
\end{lem}

\begin{proof}
  Suppose $\Gamma$ implicitly defines~$R$. Let $R'$, $\Gamma'$ be as
  above and $\Lang L' = \Lang L(\Gamma') = \Lang L \setminus \{R\}
  \cup \{R'\}$. By hypothesis, we have
  \begin{align*}
    \Gamma, \Gamma' & 
    \Sequent \lforall[x_1][\dots\lforall[x_n][(R(x_1,\dots,x_n) 
    \liff R'(x_1, \dots, x_n))]]
  \intertext{Let $c_1$, \dots,~$c_n$ be \olup{constant}s not occurring in~$\Gamma$. By the invertibility lemma we have}
    \Gamma, \Gamma', R(c_1,\dots,c_n) & \Sequent  R'(c_1, \dots, c_n)
  \intertext{Apply Maehara's Lemma to}
    \maeh{\Gamma, R(c_1,\dots,c_n)}{\Gamma'} & \Sequent
      \maeh{\ }{R'(c_1, \dots, c_n)}
  \intertext{to find an~$A \ident A(x_1, \dots, x_n)$ such that}
    \Gamma, R(c_1,\dots,c_n) & \Sequent A(c_1, \dots, c_n) \\
    A(c_1, \dots, c_n), \Gamma' & \Sequent R'(c_1, \dots, c_n)
  \intertext{have \olup{proof}s. Since $\Lang L(A) = \Lang L_1 \cap
  \Lang L_2 = \Lang L(\Gamma) \setminus \{R\}$, $A$~does not contain~$R$.
  Replace $R'$ everywhere by~$R$ in the \olup{proof} of the second sequent
  to get \olup{proof} of}
    A(c_1, \dots, c_n), \Gamma & \Sequent R(c_1, \dots, c_n)
  \end{align*}
  Apply \RightR{\lif} and \RightR{\lforall} to get \olup{proof} of
  \[
    \Gamma \Sequent \lforall[x_1][\dots\lforall[x_n][(R(x_1,\dots,x_n) \liff A(x_1, \dots, x_n))]].
  \]
  This shows that $A(x_1, \dots, x_n)$ explicitly defines $R$
  in~$\Gamma$.
\end{proof}

A third result which is equivalent to interpolation (and, like it, can
be proved using Maehara's Lemma) is the following theorem: If two
consistent theories~$\Gamma_1$, $\Gamma_2$ contain a complete theory
$\Gamma$ in the language~$\Lang L(\Gamma) \subseteq \Lang L(\Gamma_1)
\cap \Lang L(\Gamma_2)$, then $\Gamma_1 \cup \Gamma_2$ is consistent.

Recall that a complete theory is one such that for every
\olup{sentence}~$A$ in its language either $\Gamma \Proves A$ or
$\Gamma \Proves \lnot A$. Since $\Gamma_1$ and~$\Gamma_2$ are consistent
extensions of~$\Gamma$, $\Gamma$ must be consistent. It follows that
if $A$ is \olup{sentence} in the language $\Lang L(\Gamma_1) \cap \Lang
L(\Gamma_2)$, we can't have $\Gamma_1 \Entails A$ and $\Gamma_2
\Entails \lnot A$. To see this, suppose there were such an~$A$. If
$\Gamma_1 \Entails A$ then $\Gamma \Entails A$. Otherwise, by
completeness of~$\Gamma$, $\Gamma \Entails \lnot A$, and since
$\Gamma \subseteq \Gamma_1$, $\Gamma_1 \Entails \lnot A$, making
$\Gamma_1$ inconsistent. So $\Gamma \Entails A$ and since $\Gamma
\subseteq \Gamma_2$, $\Gamma_2 \Entails A$. So $\Gamma \Entails/
\lnot A$, or $\Gamma_2$ would be inconsistent.

The non-existence of \olup{sentence}~$A$ in the joint
\olup{language}~$\Lang L(\Gamma_1) \cap \Lang L(\Gamma_2)$ with $\Gamma_1
\Entails A$ and $\Gamma_2 \Entails \lnot A$ is all that's needed for
the proof, which we give here proof-theoretically using Maehara's
Lemma.

\begin{thm}[Robinson's Joint Consistency Theorem]
Suppose $\Gamma_1$, $\Gamma_2$ are consistent theories and for no~$A$
in the \olup{language} $\Lang L(\Gamma_1) \cap \Lang L(\Gamma_2)$ do we
have $\Gamma_1 \Entails A$ and $\Gamma_2 \Entails \lnot A$. Then
$\Gamma_1 \cup \Gamma_2$ is consistent.
\end{thm}

\begin{proof}
  Suppose for the contrary that $\Gamma_1 \cup \Gamma_2$ is
  inconsistent. Then the sequent $\Gamma_1, \Gamma_2 \Sequent \ $ has
  \olup{proof}. Apply Maehara's Lemma to $\maeh{\Gamma_1}{\Gamma_2}
  \Sequent \maeh{\ }{\ }$ to get \olup{sentence}~$A$ in the
  \olup{language}~$\Lang L(\Gamma_1) \cap \Lang L(\Gamma_2)$ such that
  $\Proves \Gamma_1 \Sequent A$ and $\Proves A, \Gamma_2 \Sequent
  \quad$. From the latter we get $\Proves \Gamma_2 \Sequent \lnot A$.
  But we assumed that no such \olup{sentence} exists.
\end{proof}

Above we have tacitly assumed that the theories $\Gamma$, $\Gamma_1$,
$\Gamma_2$ are finite. By compactness, the results also hold for
infinite theories.


% ==== OLP file: content/proof-theory/cut-elimination/introduction.tex ====
% Part: proof-theory
% Chapter: cut-elimination
% Section: introduction



\olfileid{pt}{cut}{int}

\olsection{Introduction}

The \Cut{} rule is the rule
\[
\Axiom$\Gamma \fCenter \Delta, A$
\Axiom$A, \Pi \fCenter \Lambda$
\RightLabel{\Cut}
\BinaryInf$\Gamma, \Pi \fCenter \Delta, \Lambda$
\DisplayProof
\]
(The \olup{formula}~$A$ is called the \emph{cut formula}.)

\Cut{} was included in Gentzen's original sequent calculus~\Log{LK}.
Gentzen's \emph{Hauptsatz} (``main theorem'') states that every
sequent \olup{provable} in~\Log{LK} is also \olup{provable} without using
the \Cut{} rule, i.e., that \Cut{} can be ``eliminated'' from
\Log{LK}-proofs. Our systems \Log{G1c} and \Log{G3c} do not contain
\Cut, but the question can be raised for them as well: can \Cut{} be
eliminated from \olup{proof}s that use \Cut{} plus the other rules of
\Log{G1c} (or those of \Log{G3c})? In our established terminology it
can be asked equivalently as: Is \Cut{} an admissible rule of
\Log{G1c} (or \Log{G3c}).

The cut elimination theorem is perhaps the most important and
celebrated result of proof theory. It first of all establishes that
\Log{G1c} and \Log{G3c} don't need a rule which you might, at first
blush, think is clearly required. For when you try to formalize actual
proofs, you'll notice that you need a way to deal with the use of
lemmas. Mathematicians like to prove results, and then appeal to these
results in proofs of other results. So you might first formalize a
proof of result $L$ (the lemma) as the \olup{proof} of a sequent
$\Sequent L$. Then you formalize the proof of the second result~$R$,
the proof of which uses~$L$, as the \olup{proof} of a sequent $L \Sequent
R$. How do you know that there is a proof of~$R$, i.e., \olup{proof} of
just $\Sequent R$? You need a way to combine the two \olup{proof}s, and
the \Cut{} rule would allow you to do that.

We mentioned that \Log{G1c} and \Log{G3c} are complete, and so they
prove any sequent you might expect to be able to prove---all valid
sequents. This can be proved directly. But at the time Gentzen was
writing, only axiomatic derivation systems were known to be complete.
To show \Log{LK} complete, Gentzen provided a translation of
\olup{derivation}s in the axiomatic system into \Log{LK}-\olup{proof}s. This
translation used \Cut{} essentially. Cut elimination is not just an
important result for classical logic: many other logics have sequent
calculi. For some logics, a direct proof of completeness for their
sequent calculi is hard or impossible, so translations from other
systems (using \Cut) are the only way to show completeness. Then
proving cut elimination proof-theoretically is necessary to establish
that \Cut{} is superfluous and that the system without \Cut{} is
complete.

Cut elimination is also important because it can be applied to obtain
other, independently interesting results. Two cases we'll consider are
the Craig interpolation lemma and Herbrand' theorem.

The ideas used to prove cut elimination typically involve showing how
\olup{proof} that uses \Cut{} can be transformed into one without. These
transformations are usually ``local,'' i.e., we take a part of !!
a{proof} (usually a sub-\olup{proof} ending in a \Cut{} inference) and
replace it with another sub-\olup{proof} of the same sequent but in which
the \Cut{} inference has been removed or replaced by others. Such
arguments provide step-by-step algorithms for transforming \olup{proof}s
with \Cut{} into ones without. This provides more information when cut
elimination procedures are used in applications. For instance, there
are model theoretic proofs of the interpolation lemma the form ``if
$A \lif B$'' is valid then an interpolant exists (never mind what
and interpolant is for the moment). A proof-theoretic argument using
cut elimination provides an algorithm that takes \olup{proof} of $A
\lif B$ and returns an interpolant. In contrast to the model
theoretic argument, the proof theoretic argument is
\emph{constructive}.

There are two popular ways to prove the cut elimination theorem. One
(going back to Gentzen) shows that you can remove top-most cuts in a
proof, and then just removes top-most cuts until none are left. The
other (originally due to Sch\"utte and Tait) focuses on the most
complex cuts in \olup{proof}, and successively removes cuts that are
maximal in the sense that no other cut has a cut \olup{formula} of higher
depth. They each have their advantages and drawbacks. For some
systems, one approach works while another is hard to implement or even
impossible. So, we describe both approaches. We'll prove cut
elimination for~\Log{G3c}. In fact, we'll show not that \Cut{} can be
eliminated, but \emph{context-sharing cut}~\CutCS:
\[
\Axiom$\Gamma \fCenter \Delta, A$
\Axiom$A, \Gamma \fCenter \Delta$
\RightLabel{\Cut}
\BinaryInf$\Gamma \fCenter \Delta$
\DisplayProof
\]
If a sequent calculus allows weakening and contraction (either as
actual rules like in \Log{G1c} or as admissible rules like in
\Log{G3c}), \Cut{} can simulate \CutCS, and vice versa. We choose
\CutCS{} over \Cut because the first cut elimination strategy is
easier to describe for~\CutCS, and the second strategy only works
for~\CutCS.


% ==== OLP file: content/proof-theory/cut-elimination/intuitionistic.tex ====
If $A \ident (B \lif C)$ then $\pi$ ends in
\[
\AxiomC{}
\RightLabel{$\pi_1'$}
\Deduce$B, \Gamma \fCenter \Delta, C$
\RightLabel{\RightR{\lif}}
\UnaryInf$\Gamma \fCenter \Delta, B \lif C$
\LeftSubproofLabel{$\pi_1$}
\AxiomC{}
\RightLabel{$\pi_2'$}
\Deduce$B \lif C, \Pi \fCenter \Lambda, B$
\AxiomC{}
\RightLabel{$\pi_2''$}
\Deduce$C, \Pi \fCenter \Lambda$
\RightLabel{\LeftR{\lif}}
\BinaryInf$B \lif C, \Pi \fCenter \Lambda$
\RightSubproofLabel{$\pi_2$}
\RightLabel{\Cut}
\BinaryInf$\Gamma, \Pi \fCenter \Delta, \Lambda$
\DisplayProof
\]
The \olup{proof} ending in \Cut,
\[
\AxiomC{}
\RightLabel{$\pi_1'$}
\Deduce$B, \Gamma \fCenter \Delta, C$
\RightLabel{\RightR{\lif}}
\UnaryInf$\Gamma \fCenter \Delta, B \lif C$
\LeftSubproofLabel{$\pi_1$}
\AxiomC{}
\RightLabel{$\pi_2'$}
\Deduce$B \lif C, \Pi \fCenter \Lambda, B$
\RightLabel{\Cut}
\BinaryInf$\Gamma, \Pi \fCenter \Delta, \Lambda, B$
\DisplayProof
\]
has lower cut \olup{height} than~$\pi$. The inductive hypthesis yields
\olup{proof}~$\pi_1''$ of $\Gamma, \Pi \fCenter \Delta, \Lambda, B$.
The \olup{proof} ending in \Cut,
\[
\AxiomC{}
\RightLabel{$\pi_1'$}
\Deduce$B!, \Gamma \fCenter \Delta, C$
\AxiomC{}
\RightLabel{$\pi_2''$}
\Deduce$C, \Pi \fCenter \Lambda$
\RightLabel{\Cut}
\BinaryInf$B, \Gamma, \Pi \fCenter \Delta, \Lambda$
\DisplayProof
\]
Both its cut rank and cut \olup{height} are less than the cut rank and
\olup{height} of~$\pi$, so the induction hypothesis yields a
\Log{G3c}-\olup{proof}~$\pi_2'''$ of $B, \Gamma, \Pi, \Pi \fCenter \Delta, \Lambda,
\Lambda$. Now apply the inductive hypothesis to the \olup{proof} ending
in \Cut,
\[
\AxiomC{}
\RightLabel{$\pi_1''$}
\Deduce$\Gamma, \Pi \fCenter \Delta, \Lambda, B$
\AxiomC{}
\RightLabel{$\pi_2'''$}
\Deduce$B, \Gamma, \Pi, \Pi \fCenter \Delta, \Lambda,
\Lambda$
\RightLabel{\Cut}
\BinaryInf$B, \Gamma, \Pi \fCenter \Delta, \Lambda$
\DisplayProof
\]
(it has cut rank $< \cutr{\pi}$) to obtain \olup{proof} of $\Gamma,
\Gamma, \Pi, \Pi, \Pi \fCenter \Delta, \Lambda, \Lambda, \Lambda$. We
obtain the \olup{proof} $\pi'$ of $\Gamma, \Pi \fCenter \Delta, \Lambda$
by applying \olref[seq][inv]{prop:G3c-cont-adm}.
% ==== OLP file: content/proof-theory/cut-elimination/midsequent.tex ====
% Part: proof-theory
% Chapter: cut-elimination
% Section: midsequent



\olfileid{pt}{cut}{mid}

\olsection{The Midsequent Theorem and Herbrand's Theorem}

An important consequence of the cut-elimination theorem is the
midsequent theorem. It states that if there is \olup{proof}~$\pi$
of~$\Gamma \Sequent \Delta$ in which every \olup{formula} is prenex, then
there is \olup{proof}~$\pi'$ containing a sequent $S \ident \Pi \Sequent
\Lambda$ (called the \emph{midsequent}) such that all inferences above
$S$ in~$\pi'$ are propositional and all inferences below~$S$ are
quantifier inferences. (!!^a{formula}~$A$ is \emph{prenex} (or in
\emph{prenex form}) if it is of the form
\[\mathsf{Q}_1x_1\dots\mathsf{Q}_nx_n\,B(x_1,\dots,x_n)\] where each
$\mathsf{Q}_i$ is either $\lforall$ or~$\lexists$.) An important
consequence of the midsequent theorem is Herbrand's theorem. Its
general case is a little hard to describe, but very generally it takes
the form: for every \olup{provable} \olup{sentence}~$A$ of first-order
logic, there is a propositional tautology~$A'$ from which $A$ can be \olup{prove}d. Here is the version for
\olup{sentence}s of the form~$\lexists[x][B(x)]$ where $B$ is
quantifier-free:

\begin{thm}[Herbrand's Theorem]
Suppose that $B(x)$ is quantifier-free and that $\Proves 
\lexists[x][B(x)]$. Then there are terms~$t_1$, \dots, $t_n$
such that $\Proves B(t_1) \lor \dots \lor B(t_n)$.
\end{thm}

The \olup{formula}~$B(t_1) \lor \dots \lor B(t_n)$ is called a
\emph{Herbrand disjunction} of~$\lexists[x][B(x)]$. Since $B$ is
quantifier-free, this disjunction is also quantifier-free. So if it is
\olup{provable}, it is \olup{provable} by a cut-free proof and consequently
with \olup{proof} that only uses propositional inferences. Hence, it is
a tautology.

The difficult part of Herbrand's Theorem is showing that a Herbrand
disjunction exists for every \olup{provable}
\olup{formula}~$\lexists[x][B(x)]$. The converse, i.e., that if
$\lexists[x][B(x)]$ has a Herbrand disjunction, then it is provable,
is easy. In fact, the Herbrand disjunction entails $\lexists[x][B(x)]$.
E.g., here is \olup{proof} in~$\Log{G3c}$ for the case where $n=2$:
\[
\Axiom$B(t_1) \fCenter \lexists[x][B(x)], B(t_1), B(t_2)$
\Axiom$B(t_2) \fCenter \lexists[x][B(x)], B(t_1), B(t_2)$
\RightLabel{\LeftR{\lor}}
\BinaryInf$B(t_1) \lor B(t_2) \fCenter \lexists[x][B(x)], B(t_1), B(t_2)$
\RightLabel{\RightR{\lexists}}
\UnaryInf$B(t_1) \lor B(t_2) \fCenter \lexists[x][B(x)], B(t_2)$
\RightLabel{\RightR{\lexists}}
\UnaryInf$B(t_1) \lor B(t_2) \fCenter \lexists[x][B(x)]$
\DisplayProof
\]

To extract the Herbrand disjunction from \olup{proof}~$\pi$ of~$\Sequent
\lexists[x][B(x)]$, consider the case where all $\RightR{\lexists}$
inferences in a cut-free \olup{proof}~$\pi$ are at the end:
\[
\AxiomC{}
\RightLabel{$\pi_1$}
\Deduce$ \fCenter \lexists[x][B(x)], B(t_1), \dots, B(t_n)$
\RightLabel{\RightR{\lexists}}
\UnaryInf$ \fCenter \lexists[x][B(x)], B(t_2), \dots, B(t_n)$
\RightLabel{$(n-1) \times \RightR{\lexists}$}
\Deduce$ \fCenter \lexists[x][B(x)]$
\DisplayProof
\]
If these are all the $\RightR{\lexists}$ inferences in~$\pi$ with
$\lexists[x][B(x)]$ active, then there are no other quantifier
inferences in~$\pi_1'$. That's because $B(x)$ contains no other
quantifiers, and the end-sequent doesn't contain any \olup{formula}s on
the left. In other words, $ \fCenter \lexists[x][B(x)], B(t_1),
\dots, B(t_n)$ is the mid-sequent of~$\pi$. Moreover, since there are
no quantifier inferences above the mid-sequent, all sequents above it
contain $\lexists[x][B(x)]$ and these are neither active nor
principal. So they can be deleted from the \olup{proof}~$\pi_1$, and we
get \olup{proof} of~$\fCenter B(t_1), \dots, B(t_n)$, and using $n-1$
applications of $\RightR{\lor}$, \olup{proof} of the Herbrand
disjunction~$\Sequent B(t_1) \lor \dots \lor B(t_n)$.

The trouble is that we can't assume that even a cut-free \olup{proof} of
$\Sequent \lexists[x][B(x)]$ has all its \RightR{\lexists} inferences
in one block at the end. There might be inferences in between these
\RightR{\lexists} inferences in which some $B(t_i)$ is principal.
Hence, the need for the midsequent theorem.

\begin{thm}[Midsequent Theorem]\ollabel{thm:midsequent} If there is a
  \Log{G3c}-\olup{proof}~$\pi$ of $\Gamma \Sequent \Delta$ where all
  \olup{formula}s in $\Gamma$ and $\Delta$ are prenex, then there is
  \olup{proof}~$\pi'$ containing a sequent $\Pi \Sequent \Lambda$ such
  that all inferences below it are quantifier inferences and all
  inferences above it are propositional inferences.
\end{thm}

\begin{proof}
  Define the \emph{order} of a quantifier inference in~$\pi$ as the
  number of propositional inferences below it, and the order~$o(\pi)$
  of $\pi$ as the sum of the orders of all its quantifier inferences.
  If $o(\pi) = 0$, then no quantifier inference has a propositional
  inference below it, and we are done: Since all quantifier inferences
  have one premise, there is a single topmost quantifier inference.
  Its premise is the midsequent.

  If $o(\pi) > 0$, pick a bottom-most quantifier inference with order
  $>0$. Since it has order $>0$, there must be a propositional
  inference below it. Since it is a bottom-most one, its conclusion
  can't be the premise of a quantifier inference: that second, lower
  quantifier inference would also have a propositional inference below
  it otherwise. We distinguish (many!) cases according to what quantifier
  inference we have, and what the propositional inference below it is.
  Because the end-sequent consists of prenex formulas only, the
  principal formula of the quantifier inference can't be active in the
  following propositional inference. Otherwise the principal formula
  of that inference would have a quantifier in the scope of a
  propositional connective. By the sub-\olup{formula} property, this would
  have to be a sub-\olup{formula} of \olup{formula} in the end-sequent.
  So, the principal \olup{formula} of the quantifier inference is
  \olup{element} of the context of the lower propositional inference.

  Here are two examples:

First, suppose we have a $\RightR{\lforall}$ which ends in the right
premise of $\LeftR{\lif}$. Then the relevant sub-\olup{proof} of~$\pi$
ends in the sub-\olup{proof}~$\pi_1$.
\[
\AxiomC{}
\RightLabel{$\pi_2$}
\Deduce$\Gamma' \fCenter \Delta', \lforall[x][A(x)], B$
\AxiomC{}
\RightLabel{$\pi_3$}
\Deduce$C, \Gamma' \fCenter \Delta', A(c)$
\RightLabel{\RightR{\lforall}}
\UnaryInf$C, \Gamma' \fCenter \Delta', \lforall[x][A(x)]$
\RightLabel{\LeftR{\lif}}
\BinaryInf$B \lif C, \Gamma' \fCenter \Delta', \lforall[x][A(x)]$
\DisplayProof
\]
We can assume that $\pi$ is regular, and so the eigenvariable~$c$ does
not occur outside~$\pi_3$; in particular, it does not occur in $B
\lif C$ or in~$\pi_2$. Now consider the proof~$\pi_1'$:
\[
\AxiomC{}
\RightLabel{$\pi_2'$}
\Deduce$\Gamma' \fCenter \Delta', \lforall[x][A(x)], A(c), B$
\AxiomC{}
\RightLabel{$\pi_3$}
\Deduce$C, \Gamma' \fCenter \Delta', \lforall[x][A(x)], A(c)$
\RightLabel{\LeftR{\lif}}
\BinaryInf$B \lif C, \Gamma' \fCenter \Delta', \lforall[x][A(x)], A(c)$
\RightLabel{\RightR{\lforall}}
\UnaryInf$B \lif C, \Gamma' \fCenter \Delta', \lforall[x][A(x)], \lforall[x][A(x)]$
\DisplayProof
\]
where $\pi_2'$ is obtained from~$\pi_2$ by weakening with~$A(c)$ on
the right. (This does not add any inferences, in particular, no
quantifier inferences which might affect the order. It also does not
violate any eigenvariable conditions in~$\pi_2$, since no constants
in~$A(c)$ are used as eigenvariables in~$\pi_2$.) The eigenvariable
condition on \RightR{\lforall} is still satisfied, since $c$ does not
occur in~$B \lif C$.

We invoke the admissibility of contraction to
obtain \olup{proof}~$\pi_1''$ of $B \lif C, \Gamma' \fCenter \Delta',
\lforall[x][A(x)]$. The proof of the admissibility of
$\RightR{\Contraction}$ for $\lforall[x][A(x)]$ and the proof of
invertibility of $\RightR{\lforall}$ used in that proof did not change
the order of any inference: it, at most, deleted inferences. Thus,
there are no more quantifier inferences in~$\pi_1''$ than in~$\pi_1$,
and each quantifier inference in~$\pi_1''$  has the same number of
proposional inferences below it as its corresponding inference
in~$\pi_1$---with the exception of the last \RightR{\lforall}: it had
a \LeftR{\lif} below it in~$\pi$, which we have moved above it
in~$\pi_1''$. Replace $\pi_1$ in $\pi$ by $\pi_1''$. The resulting
\olup{proof} has a lower order, since at least the order of the
$\RightR{\lforall}$ inference has been reduced by (at least)~$1$.
Now apply the inductive hypothesis.

The cases where the quantifier inference is \RightR{\lexists} is even
easier, since contraction is not needed and no eigenvariable condition
needs to be worried about. E.g., suppose $\RightR{\lexists}$ is
followed by $\RightR{\land}$ (this time as the left premise). The
relevant sub-\olup{proof} is $\pi_1$:
\[
\AxiomC{}
\RightLabel{$\pi_2$}
\Deduce$\Gamma' \fCenter \Delta', \lexists[x][A(x)], A(t), B$
\RightLabel{\RightR{\lexists}}
\UnaryInf$!\Gamma' \fCenter \Delta', \lexists[x][A(x)], B$
\AxiomC{}
\RightLabel{$\pi_3$}
\Deduce$\Gamma' \fCenter \Delta', \lexists[x][A(x)], C$
\RightLabel{\RightR{\land}}
\BinaryInf$\Gamma' \fCenter \Delta', \lexists[x][A(x)], B \land C$
\DisplayProof
\]
Replace $\pi_1$ in~$\pi$ by $\pi_1'$:
\[
\AxiomC{}
\RightLabel{$\pi_2$}
\Deduce$\Gamma' \fCenter \Delta', \lexists[x][A(x)], A(t), B$
\AxiomC{}
\RightLabel{$\pi_3'$}
\Deduce$\Gamma' \fCenter \Delta', \lexists[x][A(x)], A(t), C$
\RightLabel{\RightR{\land}}
\BinaryInf$\Gamma' \fCenter \Delta', \lexists[x][A(x)], A(t), B \land C$
\RightLabel{\RightR{\lexists}}
\UnaryInf$\Gamma' \fCenter \Delta', \lexists[x][A(x)], B \land C$
\DisplayProof
\]
where $\pi_3'$ is obtained from $\pi_3$ by weakening with~$A(t)$ on
the right. This weakening involves merely adding $A(t)$ to the right
side of every sequent in~$\pi_3$, so does not change the order of any
quantifier inference. The order of quantifier inferences is the same
in~$\pi'$ as in~$\pi$, except the order of the \RightR{\lexists}
inference permuted with~\RightR{\land} has decreased by~$1$. The
inductive hypothesis applies.
\end{proof}

\begin{prob}
Describe the cases in the proof of \olref{thm:midsequent} where
\LeftR{\lexists} ends in the premise of~\RightR{\lif}, and where
\LeftR{\lforall} ends in the left premise of~\LeftR{\lor}.
\end{prob}

\begin{proof}[Proof of Herbrand's Theorem]
Suppose $\pi$ ends in $\Sequent \lexists[x][A(x)]$. By
\olref{thm:midsequent}, we can transform $\pi$ into \olup{proof}~$\pi'$
with a midsequent~$S$, which must be of the form \[\Sequent
\lexists[x][A(x)], A(t_1), \dots, A(t_n).\] Since all inferences
above $S$ are propositional, $\lexists[x][A(x)]$ cannot be principal
in an inference above~$S$. Since it is not atomic, it cannot be
principal in an axiom. Since $\lexists[x][A(x)]$ cannot be a proper
subformula of $A(t_1)$ (remember that $A(x)$ is quantifier-free), it
can also not be active above~$S$. Deleting $\lexists[x][A(x)]$ from
the sub-\olup{proof} ending in~$S$ therefore yields a correct \olup{proof}
of $\Sequent A(t_1), \dots, A(t_n)$, from which \olup{proof} of
$\Sequent A(t_1) \lor \dots \lor A(t_n)$ can be obtained by $n-1$
$\RightR{\lor}$ inferences.
\end{proof}

% \begin{prob}
% Find a Herbrand disjunction of $\lexists[x][\lforall[y][(A(x) \lif
% A(y))]]$ by finding a cut-free \olup{proof} of $\Sequent
% \lexists[x][\lforall[y][(A(x) \lif A(y))]]$ and applying the
% transformation in \olref{thm:midsequent} to find a midsequent (if
% necessary).
% \end{prob}

\end{document}
